\section{Theoretical results} \label{app:proofs}
\begin{restatable}[A policy is optimal iff it induces an optimal visit distribution at every state]{lem}{optShare}\label{lem:opt-pol-visit-iff}
Let $\gamma \in (0,1)$ and let $R$ be a reward function. $\pi \in \optPi$ iff $\pi$ induces an optimal visit distribution at every state. 
\end{restatable}
\begin{proof}
By definition, a policy $\pi$ is optimal iff $\pi$ induces the maximal on-policy value at each state, which is true iff $\pi$ induces an optimal visit distribution at every state (by the dual formulation of optimal value functions). 
\end{proof}

\begin{restatable}[Transition matrix induced by a policy]{definition}{transMatrix}\label{def:trans-matrix}
$\mathbf{T}^\pi$ is the transition matrix induced by policy $\pi \in \Pi$, where $\mathbf{T}^\pi\unitvec \defeq T(s,\pi(s))$. $(\mathbf{T}^\pi)^t \unitvec$ gives the probability distribution over the states visited at time step $t$, after following $\pi$ for $t$ steps from $s$.
\end{restatable} 

\begin{restatable}[Properties of visit distribution functions]{prop}{visitDistProperties}\label{prop:visit-dist-prop} 
Let $s,s'\in\St,\fpi{s} \in \F(s)$. 
\begin{enumerate}
    \item $\fpi{s}(\gamma)$ is element-wise non-negative and element-wise monotonically increasing on $\gamma\in[0,1)$. \label{item:mono-increase}
    \item $\forall \gamma \in [0,1): \lone{\fpi{s}(\gamma)}=\geom$.\label{item:lone-visit}
\end{enumerate}
\end{restatable}
\begin{proof}
\Cref{item:mono-increase}: by examination of \cref{def:visit}, $\fpi{s} = \sum_{t=0}^\infty \prn{\gamma\mathbf{T}^\pi}^t\unitvec$. Since each $\prn{\mathbf{T}^\pi}^t$ is left stochastic and $\unitvec$ is the standard unit vector, each entry in each summand is non-negative. Therefore, $\forall \gamma \in [0,1): \fpi{s}(\gamma)^\top \unitvec[s']\geq 0$, and this function monotonically increases on $\gamma$.

\Cref{item:lone-visit}:
\begin{align}
    \lone{\fpi{s}(\gamma)}&=\lone{\sum_{t=0}^\infty \prn{\gamma\mathbf{T}^\pi}^t\unitvec}\\
    &= \sum_{t=0}^\infty\gamma^t\lone{ \prn{\mathbf{T}^\pi}^t\unitvec}\label{eq:non-negative-norm}\\
    &= \sum_{t=0}^\infty \gamma^t\label{eq:left-stoch}\\
    &= \geom.
\end{align}

\Cref{eq:non-negative-norm} follows because all entries in each $\prn{\mathbf{T}^\pi}^t\unitvec$ are non-negative by  \cref{item:mono-increase}. \Cref{eq:left-stoch} follows because each $\prn{\mathbf{T}^\pi}^t$ is left stochastic and $\unitvec$ is a stochastic vector, and so $\lone{\prn{\mathbf{T}^\pi}^t\unitvec}=1$. 
\end{proof}

\begin{restatable}[$\f\in\F(s)$ is multivariate rational on $\gamma$]{lem}{fRat}\label{f-rat}
$\fpi{}\in\F(s)$ is a multivariate rational function on $\gamma\in[0,1)$.
\end{restatable}
\begin{proof}
Let $\rf\in\rewardVS$ and consider $\fpi{}\in\F(s)$. Let $\mathbf{v}^\pi_R$ be the $\OptVf{s,\gamma}$ function in column vector form, with one entry per state value.

By the Bellman equations, $\mathbf{v}^\pi_R = \prn{\mathbf{I}-\gamma\mathbf{T}^\pi}\inv\rf.$ Let $\mathbf{A}_\gamma \defeq \prn{\mathbf{I}-\gamma\mathbf{T}^\pi}\inv$, and for state $s$, form $ \mathbf{A}_{s,\gamma}$ by replacing  $\mathbf{A}_\gamma$'s column for state $s$ with $\rf$. As noted by \citet{lippman1968set}, by Cramer's rule, $V^\pi_R(s,\gamma)=\frac{\det{\mathbf{A}_{s,\gamma}}}{\det\mathbf{A}_\gamma}$ is a rational function with numerator and denominator having degree at most $\abs{\St}$.

In particular, for each state indicator reward function $\unitvec[s_i]$, $V^\pi_{s_i}(s,\gamma)=\fpi{s}(\gamma)^\top\unitvec[s_i]$ is a rational function of $\gamma$ whose numerator and denominator each have degree at most $\abs{\St}$. This implies that $\f^\pi(\gamma)$ is multivariate rational on $\gamma\in[0,1)$.
\end{proof}

\begin{restatable}[On-policy value is rational on $\gamma$]{cor}{smoothOnPol}\label{smoothOnPol}
Let $\pi\in\Pi$ and $R$ be any reward function. $V^\pi_R(s,\gamma)$ is rational on  $\gamma\in[0,1)$. 
\end{restatable}
\begin{proof}
$V^\pi_R(s,\gamma)=\fpi{s}(\gamma)^\top \rf$, and $\f$ is a multivariate rational function of $\gamma$ by \cref{f-rat}. Therefore, for fixed $\rf$, $\fpi{s}(\gamma)^\top \rf$ is a rational function of $\gamma$.
\end{proof}

\subsection{Non-dominated visit distribution functions}\label{sec:nondom}

\begin{restatable}[Continuous reward function distribution]{definition}{dCont}
Results with $\Dcont$ hold for any absolutely continuous reward function distribution.
\end{restatable}

\begin{remark}
We assume $\rewardVS$ is endowed with the standard topology.
\end{remark}

\begin{restatable}[Distinct linear functionals disagree almost everywhere on their domains]{lem}{distinctLin}\label{lem:distinct-lin-prob}
Let $\x,\x'\in\rewardVS$ be distinct. $\optprob[\rf\sim\Dcont]{\x^\top \rf = \x'^\top\rf}=0$.
\end{restatable}
\begin{proof}
$\set{\rf \in\rewardVS\mid (\x-\x')^\top \rf=0}$ is a hyperplane since $\x-\x'\neq \mathbf{0}$. Therefore, it has no interior in the standard topology on $\rewardVS$. Since this empty-interior set is also convex, it has zero Lebesgue measure. By the Radon-Nikodym theorem, it has zero measure under any continuous distribution $\Dcont$.
\end{proof} 

\begin{restatable}[Unique maximization of almost all vectors]{cor}{uniqueMax}\label{cor:distinct-maximized}
Let $X\subsetneq\rewardVS$ be finite. $\optprob[\rf\sim\Dcont]{\abs{\argmax_{\x''\in X} \x''^\top \rf}>1}=0$. 
\end{restatable}
\begin{proof}
Let $\x,\x'\in X$ be distinct. For any $\rf\in\rewardVS$, $\x,\x'\in \argmax_{\x''\in X} \x''^\top \rf$ iff $\x^\top \rf = \x'^\top\rf\geq \max_{\x''\in X\setminus \set{\x,\x'}} \x''^\top \rf$. By \cref{lem:distinct-lin-prob}, $\x^\top \rf = \x'^\top\rf$ holds with probability 0 under any $\Dcont$.
\end{proof} 

\subsubsection{Generalized non-domination results}
Our formalism includes both $\Fnd(s)$ and $\RSDnd$; we therefore prove results that are applicable to both. 

\begin{restatable}[Non-dominated linear functionals]{definition}{ndLinFunc}\label{def:nd-lin-func}
Let $X\subsetneq \rewardVS$ be finite. $\ND{X}\defeq \set{\x\in X\mid \exists \rf \in \rewardVS: \x^\top \rf > \max_{\x'\in X\setminus \set{\x}} \x'^\top \rf}$.
\end{restatable}

\begin{restatable}[All vectors are maximized by a non-dominated linear functional]{lem}{allRFMaxND}\label{lem:all-rf-max-nd}
Let $\rf \in \rewardVS$ and let $X\subsetneq \rewardVS$ be finite and non-empty. $\exists \x^* \in \ND{X}: \x^{*\top}\rf=\max_{\x \in X} \x^\top \rf$.
\end{restatable}
\begin{proof}
Let $A(\rf\mid X)\defeq \argmax_{\x\in X} \x^\top \rf=\set{\x_1,\ldots, \x_n}$. Then 
\begin{align}
    \x_1^\top \rf = \cdots= \x_n^\top \rf> \max_{\x' \in X \setminus A(\rf\mid X)} \x'^\top \rf.\label{eq:expr-cont-nd-func}
\end{align}

In \cref{eq:expr-cont-nd-func}, each $\x^\top \rf$ expression is linear on $\rf$. The $\max$ is piecewise linear on $\rf$ since it is the maximum of a finite set of linear functionals. In particular, all expressions in \cref{eq:expr-cont-nd-func} are continuous on $\rf$, and so we can find some $\delta>0$ neighborhood $B(\rf,\delta)$ such that $\forall \rf' \in B(\rf,\delta): \max_{\x_i\in A(\rf\mid X)} \x_i^\top \rf' > \max_{\x'\in X\setminus A(\rf\mid X)} \x'^\top \rf'$.

But almost all $\rf' \in B(\rf,\delta)$ are maximized by a unique functional $\x^*$ by \cref{cor:distinct-maximized}; in particular, at least one such $\rf''$ exists. Formally, $\exists \rf''\in B(\rf,\delta):\x^{*\top} \rf''>\max_{\x' \in X\setminus \set{\x^*}} \x'^\top \rf''$. Therefore, $\x^*\in \ND{X}$ by \cref{def:nd-lin-func}. 

$\x^{*\top} \rf'\geq\max_{\x_i\in A(\rf\mid X)} \x_i^\top \rf' > \max_{\x'\in X\setminus A(\rf\mid X)} \x'^\top \rf'$, with the strict inequality following because  $\rf'' \in B(\rf,\delta)$. These inequalities imply that $\x^*\in A(\rf\mid X)$.
\end{proof}

\begin{restatable}[Maximal value is invariant to restriction to non-dominated functionals]{cor}{ndFuncIndif}\label{cor:nd-func-indif}
Let $\rf \in \rewardVS$ and let $X\subsetneq \rewardVS$ be finite. $\max_{\x\in X} \x^\top \rf=\max_{\x\in \ND{X}} \x^\top \rf$. 
\end{restatable}
\begin{proof}
If $X$ is empty, holds trivially. Otherwise, apply \cref{lem:all-rf-max-nd}.
\end{proof}

\begin{restatable}[How non-domination containment affects optimal value]{lem}{ndOptContain}\label{lem:nd-opt-contain}
Let $\rf \in \rewardVS$ and let $X,X'\subsetneq \rewardVS$ be finite. 
\begin{enumerate}
    \item If $\ND{X}\subseteq X'$, then $\max_{\x\in X} \x^\top \rf\leq\max_{\x'\in X'} \x'^\top \rf$.\label{item:ND-contain-max}
    \item If $\ND{X}\subseteq X'\subseteq X$, then $\max_{\x\in X} \x^\top \rf=\max_{\x'\in X'} \x'^\top \rf$.\label{item:ND-contain-max-2}
\end{enumerate}
\end{restatable}
\begin{proof}
\Cref{item:ND-contain-max}: 
\begin{align}
\max_{\x\in X} \x^\top \rf&=\max_{\x\in \ND{X}} \x^\top \rf\label{eq:item-max}\\
&\leq \max_{\x'\in X'} \x'^\top \rf.\label{eq:item-max-2}
\end{align}
\Cref{eq:item-max} follows by \cref{cor:nd-func-indif}. \Cref{eq:item-max-2} follows because $\ND{X}\subseteq X'$.

\Cref{item:ND-contain-max-2}: by \cref{item:ND-contain-max}, $\max_{\x\in X} \x^\top \rf\leq\max_{\x'\in X'} \x'^\top \rf$. Since $X'\subseteq X$, we also have $\max_{\x\in X} \x^\top \rf\geq\max_{\x'\in X'} \x'^\top \rf$, and so equality must hold.
\end{proof}

\begin{restatable}[Non-dominated vector functions]{definition}{ndVecFun}\label{def:nd-vec-func} 
Let $I\subseteq \reals$ and let $F\subsetneq \prn{\rewardVS}^I$ be a finite set of vector-valued functions on $I$. $\ND{F}\defeq \set{\f\in F\mid \exists \gamma\in I, \rf \in \rewardVS: \f(\gamma)^\top \rf> \max_{\f'\in F\setminus \set{\f}}\f'(\gamma)^\top \rf}$.
\end{restatable}

\begin{remark}
$\Fnd(s)=\ND{\F(s)}$ by \cref{def:nd}.
\end{remark}

\begin{restatable}[Affine transformation of visit distribution sets]{definition}{affTransfShorthand}\label{def:aff-transf-short}
For notational convenience, we define set-scalar multiplication and set-vector addition on $X\subseteq \rewardVS$: for $c\in\reals$, $cX\defeq \set{c\x\mid\x\in X}$. For $\av \in \rewardVS$, $X+\av\defeq \set{\x+\av\mid \x \in X}$. Similar operations hold when $X$ is a set of vector functions $\reals\mapsto \rewardVS$.
\end{restatable}

\begin{restatable}[Invariance of non-domination under positive affine transform]{lem}{posAffNDInvar}\label{lem:pos-aff-nd-invar}\hfill
\begin{enumerate}
    \item Let $X\subsetneq \rewardVS$ be finite. If $\x\in \ND{X}$, then $\forall c>0, \av\in\rewardVS: (c\x+\av) \in \ND{cX+\av}$.\label{item:invar-vectors}
    \item Let $I\subseteq \reals$ and let $F\subsetneq \prn{\rewardVS}^I$ be a finite set of vector-valued functions on $I$. If $\f \in \ND{F}$, then $\forall c>0, \av\in\rewardVS: (c\f+\av) \in \ND{cF+\av}$.\label{item:invar-vector-fns}
\end{enumerate}
\end{restatable}
\begin{proof}
\Cref{item:invar-vectors}: Suppose $\x\in \ND{X}$ is strictly optimal for $\rf\in\rewardVS$. Then let $c>0,\av\in\rewardVS$ be arbitrary, and define $b\defeq \av^\top \rf$. 
\begin{align}
    \x^\top\rf &> \max_{\x'\in X\setminus\set{\x}} \x'^\top \rf\\
    c\x^\top\rf+b &> \max_{\x'\in X\setminus\set{\x}} c\x'^\top \rf+b\label{eq:rescale-c}\\
    (c\x+\av)^\top\rf &> \max_{\x'\in X\setminus\set{\x}} (c\x'+\av)^\top \rf\label{eq:b-defn}\\
    (c\x+\av)^\top\rf &> \max_{\x''\in \prn{cX+\av}\setminus\set{c\x+\av}} \x''^\top \rf.
\end{align}
\Cref{eq:rescale-c} follows because $c>0$. \Cref{eq:b-defn} follows by the definition of $b$.

\Cref{item:invar-vector-fns}: If $\f\in\ND{F}$, then by \cref{def:nd-vec-func}, there exist $\gamma\in I, \rf\in\rewardVS$ such that
\begin{align}
    \f(\gamma)^\top \rf> \max_{\f'\in F\setminus \set{\f}}\f'(\gamma)^\top \rf.
\end{align}

Apply \cref{item:invar-vectors} to conclude
\begin{align}
    (c\f(\gamma)+\av)^\top \rf> \max_{(c\f'+\av)\in (cF+\av)\setminus \set{c\f+\av}}(c\f'(\gamma)+\av)^\top \rf.
\end{align}

Therefore, $(c\f+\av) \in \ND{cF+\av}$.
\end{proof}

\subsubsection{Inequalities which hold under most reward function distributions}

\ineqMost*

\begin{restatable}[Helper lemma for demonstrating $\geqMost$]{lem}{helperGeqMost}\label{lem:helper-geq-most}
Let $\distSet\subseteq \Delta(\rewardVS)$. If $\exists \phi\in\mdpPermGroup$ such that for all $\D\in\distSet$, $f_1\prn{\D}< f_2\prn{\D}$ implies that $f_1\prn{\phi\cdot \D}> f_2\prn{\phi\cdot \D}$, then $f_1(\D) \geqMost[][\distSet] f_2(\D)$.
\end{restatable}
\begin{proof}
Since $\phi$ does not belong to the stabilizer of $\mdpPermGroup$, $\phi$ acts injectively on $\orbi[\D]$. By assumption on $\phi$, the image of $\{\D'\in\orbi[\D]\mid f_1(\D')<f_2(\D')\}$ under $\phi$ is a subset of $\{\D'\in\orbi[\D]\mid f_1(\D')>f_2(\D')\}$. Since $\phi$ is injective, $\abs{\{\D'\in\orbi[\D]\mid f_1(\D')<f_2(\D')\}}\leq \abs{\{\D'\in\orbi[\D]\mid f_1(\D')>f_2(\D')\}}$. $f_1(\D) \geqMost[][\distSet] f_2(\D)$ by \cref{def:ineq-most-dists}.
\end{proof}

\begin{restatable}[A helper result for expectations of functions]{lem}{helperPerm}\label{lem:helper-perm}
Let $B_1,\ldots,B_n\subsetneq \rewardVS$ be finite and let $\distSet \subseteq \Delta(\rewardVS)$. Suppose $f$ is a function of the form
\begin{align}
    f\prn{B_1,\ldots,B_n \mid \D}=\E{\rf\sim \D}{g\prn{\max_{\bv_1\in B_1} \bv_1^\top \rf,\ldots, \max_{\bv_n\in B_n} \bv_n^\top \rf}}\label{eq:helper}
\end{align}
for some function $g$, and that $f$ is well-defined for all $\D \in \distSet$. Let $\phi$ be a state permutation. Then
\begin{equation}
    f\prn{B_1,\ldots,B_n \mid \D}=f\prn{\phi\cdot B_1 ,\ldots,\phi \cdot B_n \mid \phi \cdot\D}.
\end{equation}
\end{restatable}
\begin{proof} Let distribution $\D$ have probability measure $F$, and let $\phi\cdot \D$ have probability measure $F_\phi$. 
\begin{align}
    &f\prn{B_1,\ldots,B_n \mid \D}\\
    \defeq{}&\E{\rf\sim \D}{g\prn{\max_{\bv_1\in B_1} \bv_1^\top \rf,\ldots, \max_{\bv_n\in B_n} \bv_n^\top \rf}}\\
    \defeq{}&\int_{\rewardVS} g\prn{\max_{\bv_1\in B_1} \bv_1^\top \rf,\ldots, \max_{\bv_n\in B_n} \bv_n^\top \rf} \dF[\rf][F]\\
    ={}&\int_{\rewardVS} g\prn{\max_{\bv_1\in B_1} \bv_1^\top \rf,\ldots, \max_{\bv_n\in B_n} \bv_n^\top \rf} \dF[\permute\rf][F_\phi]\label{eq:permute-prob-meas}\\
    ={}&\int_{\rewardVS} g\prn{\max_{\bv_1\in B_1} \bv_1^\top \prn{\permute\inv\rf'},\ldots, \max_{\bv_n\in B_n} \bv_n^\top \prn{\permute\inv\rf'}} \abs{\det \permute}\dF[\rf'][F_\phi]\label{eq:change-of-variables}\\
    ={}&\int_{\rewardVS} g\prn{\max_{\bv_1\in B_1} \prn{\permute\bv_1}^\top \rf',\ldots, \max_{\bv_n\in B_n} \prn{\permute\bv_n}^\top \rf'} \dF[\rf'][F_\phi]\label{eq:change-of-variables-2}\\
    ={}&\int_{\rewardVS} g\prn{\max_{\bv_1'\in \phi\cdot B_1} \bv_1'^{\top} \rf',\ldots, \max_{\bv_n'\in \phi\cdot B_n} \bv_n'^{\top} \rf'} \dF[\rf'][F_\phi]\\
    \eqdef{}& f\prn{\phi \cdot B_1,\ldots,\phi\cdot B_n \mid \phi\cdot \D}.
\end{align}

\Cref{eq:permute-prob-meas} follows by the definition of $F_\phi$ (\cref{def:pushforward-permute}). \Cref{eq:change-of-variables} follows by substituting $\rf'\defeq \permute\rf$. \Cref{eq:change-of-variables-2} follows from the fact that all permutation matrices have unitary determinant and are orthogonal (and so $(\permute\inv)^\top=\permute$). 
\end{proof}

\begin{restatable}[Support of $\Dany$]{definition}{DefSupportDist}\label{def:support-dist}
Let $\Dany$ be any reward function distribution. $\supp[\Dany]$ is the smallest closed subset of $\rewardVS$ whose complement has measure zero under $\Dany$. 
\end{restatable} 

\begin{restatable}[Linear functional optimality probability]{definition}{linFNProb}
For finite $A,B\subsetneq \rewardVS$, the \emph{probability under $\Dany$ that $A$ is optimal over $B$} is $\phelper{A\geq B}[\Dany]\defeq \optprob[\rf \sim \Dany]{\max_{\av\in A} \av^\top \rf \geq \max_{\bv\in B} \bv^\top \rf}$. 
\end{restatable}

\begin{restatable}[Non-dominated linear functionals and their optimality probability]{prop}{helperPositiveProb}\label{prop:helper-positive-prob} 
Let $A\subsetneq \rewardVS$ be finite. If $\exists b<c:[b,c]^{\abs{\St}} \subseteq \supp[\Dany]$, then $\av\in\ND{A}$ implies that $\av$ is strictly optimal for a set of reward functions with positive measure under $\Dany$. 
\end{restatable}
\begin{proof}
Suppose $\exists b<c:[b,c]^{\abs{\St}} \subseteq \supp[\Dany]$. If $\av\in\ND{A}$, then let $\rf$ be such that $\av^\top \rf > \max_{\av'\in A\setminus\set{\av}} \av'^\top \rf$. For $a_1>0, a_2\in \reals$, positively affinely transform $\rf'\defeq a_1\rf + a_2\mathbf{1}$ (where $\mathbf{1}\in\rewardVS$ is the all-ones vector) so that $\rf' \in (b,c)^{\abs{\St}}$. 

Note that $\av$ is still strictly optimal for $\rf'$:
\begin{equation}
    \av^\top \rf > \max_{\av'\in A\setminus\set{\av}} \av'^\top \rf \iff \av^\top \rf' > \max_{\av'\in A\setminus\set{\av}} \av'^\top \rf'.\label{eq:cont-lin-ineq}
\end{equation}

Furthermore, by the continuity of both terms on the right-hand side of \cref{eq:cont-lin-ineq}, $\av$ is strictly optimal for reward functions in some open neighborhood $N$ of $\rf'$. Let $N'\defeq N\cap (b,c)^{\abs{\St}}$. $N'$ is still open in $\rewardVS$ since it is the intersection of two open sets $N$ and $(b,c)^{\abs{\St}}$. 

$\Dany$ must assign positive probability measure to all open sets in its support; otherwise, its support would exclude these zero-measure sets by \cref{def:support-dist}. Therefore, $\Dany$ assigns positive probability to $N'\subseteq \supp[\Dany]$. 
\end{proof}   

\begin{restatable}[Expected value of similar linear functional sets]{lem}{ndFuncIncr}\label{lem:nd-func-incr-pwr}
Let $A,B\subsetneq \rewardVS$ be finite, let $A'$ be such that $\ND{A}\subseteq A'\subseteq A$, and let $g:\reals\to\reals$ be an increasing function. If $B$ contains a copy $B'$ of $A'$ via $\phi$, then
\begin{equation}
    \E{\rf \sim \Dbd}{g\prn{\max_{\av\in A}\av^\top \rf}} \leq \E{\rf \sim \phi\cdot \Dbd}{g\prn{\max_{\bv\in B}\bv^\top \rf}}.\label{eq:nd-func-incr}
\end{equation} 

If $\ND{B}\setminus B'$ is empty, then \cref{eq:nd-func-incr} is an equality. If $\ND{B}\setminus B'$ is non-empty, $g$ is strictly increasing, and $\exists b<c: (b,c)^{\abs{\St}}\subseteq\supp[\Dbd]$, then \cref{eq:nd-func-incr} is strict.
\end{restatable}
\begin{proof}
Because $g:\reals\to\reals$ is increasing, it is measurable (as is $\max$). Therefore, the relevant expectations exist for all $\Dbd$.
\begin{align}
     \E{\rf \sim \Dbd}{g\prn{\max_{\av\in A}\av^\top \rf}} &=\E{\rf \sim \Dbd}{g\prn{\max_{\av\in A'}\av^\top \rf}}\label{eq:nd-func-incr-restr}\\
     &=\E{\rf \sim \phi\cdot \Dbd}{g\prn{\max_{\av\in \phi\cdot A'}\av^\top \rf}}\label{eq:apply-helper}\\
     &=\E{\rf \sim \phi\cdot \Dbd}{g\prn{\max_{\bv\in B'}\bv^\top \rf}}\label{eq:sim-helper}\\
     &\leq\E{\rf \sim \phi\cdot \Dbd}{g\prn{\max_{\bv\in B}\bv^\top \rf}}.\label{eq:sim-inequality}
\end{align}
\Cref{eq:nd-func-incr-restr} holds because $\forall \rf\in\rewardVS:\max_{\av\in A}\av^\top \rf=\max_{\av\in A'}\av^\top \rf$ by \cref{lem:nd-opt-contain}'s \cref{item:ND-contain-max-2} with $X\defeq A$, $X'\defeq A'$. \Cref{eq:apply-helper} holds by \cref{lem:helper-perm}. \Cref{eq:sim-helper} holds by the definition of $B'$. Furthermore, our assumption on $\phi$ guarantees that $B'\subseteq B$. Therefore, $\max_{\bv\in B'}\bv^\top \rf\leq \max_{\bv\in B}\bv^\top \rf$, and so \cref{eq:sim-inequality} holds by the fact that $g$ is an increasing function. Then \cref{eq:nd-func-incr} holds. 

If $\ND{B}\setminus B'$ is empty, then $\ND{B}\subseteq B'$. By assumption, $B'\subseteq B$. Then apply \cref{lem:nd-opt-contain} \cref{item:ND-contain-max-2} with $X\defeq B$, $X'\defeq B'$ in order to conclude that \cref{eq:sim-inequality} is an equality. Then \cref{eq:nd-func-incr} is also an equality.

Suppose that $g$ is strictly increasing, $\ND{B}\setminus B'$ is non-empty, and $\exists b<c: (b,c)^{\abs{\St}}\subseteq\supp[\Dbd]$. Let $\x \in \ND{B}\setminus B'$.
\begin{align}
     \E{\rf \sim \phi\cdot \Dbd}{g\prn{\max_{\bv\in B'}\bv^\top \rf}}&< \E{\rf \sim \phi\cdot \Dbd}{g\prn{\max_{\av\in B'\cup\set{\x}}\bv^\top \rf}}\label{eq:sim-helper-2}\\
     &\leq\E{\rf \sim \phi\cdot \Dbd}{g\prn{\max_{\bv\in B}\bv^\top \rf}}.\label{eq:sim-inequality-2}
\end{align}

$\x$ is strictly optimal for a positive-probability subset of $\supp[\Dbd]$ by \cref{prop:helper-positive-prob}. Since  $g$ is strictly increasing, \cref{eq:sim-helper-2} is strict. Therefore, we conclude that \cref{eq:nd-func-incr} is strict.
\end{proof}

\begin{restatable}[For continuous\ {\iid} distributions $\Diid$, $\exists b<c: (b,c)^{\abs{\St}}\subseteq \optSupp(\Diid)$]{lem}{IIDContains}\label{lem:iid-contains}
\end{restatable}
\begin{proof}
$\Diid\defeq \Dist^{\abs{\St}}$. Since the state reward distribution $\Dist$ is continuous, $\Dist$ must have support on some open interval $(b,c)$. Since $\Diid$ is {\iid} across states, $(b,c)^{\abs{\St}}\subseteq \supp[\Diid]$.
\end{proof} 

\begin{restatable}[Bounded, continuous {\iid} reward]{definition}{BCIIDDef}
$\DSetBCiid$ is the set of $\Diid$ which equal $\Dist^{\abs{\St}}$ for some continuous, bounded-support distribution $\Dist$ over $\reals$.
\end{restatable}
 
\begin{restatable}[Expectation superiority lemma]{lem}{expectSuperior}\label{lem:expect-superior}
Let $A,B\subsetneq \rewardVS$ be finite and let $g:\reals\to \reals$ be an increasing function. If $B$ contains a copy $B'$ of $\ND{A}$ via $\phi$, then 
\begin{align}
    \E{\rf \sim \Dbd}{g\prn{\max_{\av\in A}\av^\top \rf}}\leqMost[][\DSetBd]  \E{\rf \sim \Dbd}{g\prn{\max_{\bv\in B}\bv^\top \rf}}.\label{eq:permute-superior}
\end{align}

Furthermore, if $g$ is strictly increasing and $\ND{B}\setminus \phi\cdot \ND{A}$ is non-empty, then \cref{eq:permute-superior} is strict for all $\Diid\in\DSetBCiid$. In particular, $\E{\rf \sim \Dbd}{g\prn{\max_{\av\in A}\av^\top \rf}}\not\geqMost[][\DSetBd]  \E{\rf \sim \Dbd}{g\prn{\max_{\bv\in B}\bv^\top \rf}}$.
\end{restatable}
\begin{proof}
Because $g:\reals\to\reals$ is increasing, it is measurable (as is $\max$). Therefore, the relevant expectations exist for all $\Dbd$.

Suppose that $\Dbd$ is such that $\E{\rf \sim \Dbd}{g\prn{\max_{\bv\in B}\bv^\top \rf}}< \E{\rf \sim \Dbd}{g\prn{\max_{\av\in A}\av^\top \rf}}$. 
\begin{align}
    \E{\rf \sim \phi\cdot \Dbd}{g\prn{\max_{\av\in A}\av^\top \rf}}&\leq\E{\rf \sim \phi^2\cdot \Dbd}{g\prn{\max_{\bv\in B}\bv^\top \rf}}\label{eq:first-leq}\\
    &=\E{\rf \sim \Dbd}{g\prn{\max_{\bv\in B}\bv^\top \rf}}\label{eq:involve}\\
    &<\E{\rf \sim \Dbd}{g\prn{\max_{\av\in A}\av^\top \rf}}\label{eq:superior-leq}\\
    &\leq  \E{\rf \sim \phi\cdot \Dbd}{g\prn{\max_{\bv\in B}\bv^\top \rf}}.\label{eq:second-leq}
\end{align}
\Cref{eq:first-leq} follows by applying \cref{lem:nd-func-incr-pwr} with permutation $\phi$ and $A'\defeq \ND{A}$. \Cref{eq:involve} follows because involutions satisfy $\phi\inv=\phi$, and $\phi^2$ is therefore the identity. \Cref{eq:superior-leq} follows because we assumed that $\E{\rf \sim \Dbd}{g\prn{\max_{\bv\in B}\bv^\top \rf}}< \E{\rf \sim \Dbd}{g\prn{\max_{\av\in A}\av^\top \rf}}$. \Cref{eq:second-leq} follows by applying \cref{lem:nd-func-incr-pwr} with permutation $\phi$ and and $A'\defeq \ND{A}$. By \cref{lem:helper-geq-most}, \cref{eq:permute-superior} holds.

Suppose $g$ is strictly increasing and $\ND{B}\setminus B'$ is non-empty. Let $\phi'\in\mdpPermGroup$.
\begin{align}
\E{\rf \sim \phi'\cdot \Diid}{g\prn{\max_{\av\in A}\av^\top \rf}}&=\E{\rf \sim \Diid}{g\prn{\max_{\av\in A}\av^\top \rf}}\label{eq:same-iid-dist-exp}\\
&<\E{\rf \sim \phi\cdot \Diid}{g\prn{\max_{\bv\in B}\bv^\top \rf}}\label{eq:sim-iid-exp}\\
&=\E{\rf \sim \phi'\cdot \Diid}{g\prn{\max_{\bv\in B}\bv^\top \rf}}.\label{eq:same-iid-dist-exp-2}
\end{align}

\Cref{eq:same-iid-dist-exp} and \cref{eq:same-iid-dist-exp-2} hold because $\Diid$ distributes reward identically across states: $\forall \phi_x\in\mdpPermGroup: \phi_x\cdot \Diid=\Diid$. By \cref{lem:iid-contains}, $\exists b<c: (b,c)^{\abs{\St}}\subseteq \supp[\Diid]$. Therefore, apply \cref{lem:nd-func-incr-pwr} with $A'\defeq \ND{A}$ to conclude that \cref{eq:sim-iid-exp} holds.

Therefore, $\forall \phi' \in \mdpPermGroup: \E{\rf \sim \phi'\cdot \Diid}{g\prn{\max_{\av\in A}\av^\top \rf}}<\E{\rf \sim \phi'\cdot \Diid}{g\prn{\max_{\bv\in B}\bv^\top \rf}}$, and so $\E{\rf \sim \Dbd}{g\prn{\max_{\av\in A}\av^\top \rf}}\not\geqMost[][\DSetBd]  \E{\rf \sim \Dbd}{g\prn{\max_{\bv\in B}\bv^\top \rf}}$ by \cref{def:ineq-most-dists}.
\end{proof}

\begin{restatable}[Indicator function]{definition}{indicDef}
Let $L$ be a predicate which takes input $x$. $\indic{L(x)}$ is the function which returns 1 when $L(x)$ is true, and 0 otherwise.
\end{restatable} 

\begin{restatable}[Optimality probability inclusion relations]{lem}{optprbInclusion}\label{lem:inclusion-opt}
Let $X,Y\subsetneq \rewardVS$ be finite and suppose $Y'\subseteq Y$. \begin{equation}
    \phelper{X\geq Y}[\Dany]\leq\phelper{X\geq Y'}[\Dany]\leq \phelper{X\cup\prn{Y\setminus Y'}\geq Y}[\Dany].\label{eq:inclusion-opt}
\end{equation}

If $\exists b<c: (b,c)^{\abs{\St}}\subseteq \supp[\Dany]$, $X\subseteq Y$, and $\ND{Y}\cap \prn{Y\setminus Y'}$ is non-empty, then the second inequality is strict.
\end{restatable}
\begin{proof}
\begin{align}
    \phelper{X\geq Y}[\Dany]&\defeq \E{\rf \sim \Dany}{\indic{\max_{\x\in X} \x^\top \rf \geq \max_{\mathbf{y} \in Y} \mathbf{y}^\top \rf}}\\
    &\leq \E{\rf \sim \Dany}{\indic{\max_{\x\in X} \x^\top \rf \geq \max_{\mathbf{y} \in Y'} \mathbf{y}^\top \rf}}\label{eq:leq-Y-contain}\\
    &\leq \E{\rf \sim \Dany}{\indic{\max_{\x\in X\cup (Y\setminus Y')} \x^\top \rf \geq \max_{\mathbf{y} \in Y'} \mathbf{y}^\top \rf}}\label{eq:leq-Y-union}\\
    &= \E{\rf \sim \Dany}{\indic{\max_{\x\in X\cup (Y\setminus Y')} \x^\top \rf \geq \max_{\mathbf{y} \in Y'\cup (Y\setminus Y')} \mathbf{y}^\top \rf}}\label{eq:leq-Y-union-2}\\
    &=\E{\rf \sim \Dany}{\indic{\max_{\x\in X\cup (Y\setminus Y')} \x^\top \rf \geq \max_{\mathbf{y} \in Y} \mathbf{y}^\top \rf}}\\
    &\eqdef \phelper{X\cup\prn{Y\setminus Y'}\geq Y}[\Dany].
\end{align}
\Cref{eq:leq-Y-contain} follows because $\forall \rf \in \rewardVS: \indic{\max_{\x\in X} \x^\top \rf \geq \max_{\mathbf{y} \in Y} \mathbf{y}^\top \rf}\leq \indic{\max_{\x\in X} \x^\top \rf \geq \max_{\mathbf{y} \in Y'} \mathbf{y}^\top \rf}$ since $Y' \subseteq Y$; note that \cref{eq:leq-Y-contain} equals $\phelper{X\geq Y'}[\Dany]$, and so the first inequality of \cref{eq:inclusion-opt} is shown. \Cref{eq:leq-Y-union} holds because $\forall \rf \in \rewardVS: \indic{\max_{\x\in X} \x^\top \rf \geq \max_{\mathbf{y} \in Y'} \mathbf{y}^\top \rf}\leq \indic{\max_{\x\in X\cup (Y\setminus Y')} \x^\top \rf \geq \max_{\mathbf{y} \in Y'} \bv^\top \rf}$. 

Suppose $\exists b<c: (b,c)^{\abs{\St}}\subseteq \supp[\Dany]$, $X\subseteq Y$, and $\ND{Y}\cap \prn{Y\setminus Y'}$ is non-empty. Let $\mathbf{y}^*\in \ND{Y}\cap \prn{Y\setminus Y'}$. By \cref{prop:helper-positive-prob}, $\mathbf{y}^*$ is strictly optimal on a subset of $\supp[\Dany]$ with positive measure under $\Dany$. In particular, for a set of $\rf^*$ with positive measure under $\Dany$, we have $\mathbf{y}^{*\top}\rf^*> \max_{\mathbf{y}\in Y'} \mathbf{y}^\top \rf^*.$

Then \cref{eq:leq-Y-union} is strict, and therefore the second inequality of \cref{eq:inclusion-opt} is strict as well.
\end{proof}

\begin{restatable}[Optimality probability of similar linear functional sets]{lem}{optProbSim}\label{lem:sim-lin-func-opt}
Let $A,B,C\subsetneq \rewardVS$ be finite, and let $Z\subseteq \rewardVS$ be such that $\ND{C}\subseteq Z\subseteq C$. If $\ND{A}$ is similar to $B'\subseteq B$ via $\phi$ such that $\phi\cdot \prn{Z\setminus \prn{B\setminus B'}}=Z\setminus \prn{B\setminus B'}$, then 
\begin{equation}
    \phelper{A\geq C}[\Dany]\leq \phelper{B\geq C}[\phi\cdot \Dany].\label{eq:sim-nd-func}
\end{equation} 
If $B'=B$, then \cref{eq:sim-nd-func} is an equality. If $\exists b<c: (b,c)^{\abs{\St}}\subseteq \supp[\Dany]$, $B'\subseteq C$, and $\ND{C}\cap \prn{B\setminus B'}$ is non-empty, then \cref{eq:sim-nd-func} is strict.
\end{restatable}
\begin{proof}
\begin{align}
    \phelper{A\geq C}[\Dany]&=\phelper{A\geq Z}[\Dany]\label{eq:nd-restrict-opt-lin}\\ 
    &=\phelper{\ND{A}\geq Z}[\Dany]\label{eq:nd-restrict-opt-lin-2}\\
    &\leq \phelper{\ND{A}\geq Z\setminus\prn{B\setminus B'}}[\Dany]\label{eq:fewer-requirements}\\
    &= \phelper{\phi\cdot \ND{A}\geq \phi\cdot Z\setminus\prn{B\setminus B'}}[\phi\cdot \Dany]\label{eq:phi-helper-nd-opt}\\
    &= \phelper{B'\geq Z\setminus\prn{B\setminus B'}}[\phi\cdot \Dany]\label{eq:phi-assumptions}\\
    &\leq \phelper{B'\cup \prn{B\setminus B'}\geq Z}[\phi\cdot \Dany]\label{eq:reunion}\\
    &= \phelper{B\geq C}[\phi\cdot \Dany].\label{eq:final-phelper-leq}
\end{align}

\Cref{eq:nd-restrict-opt-lin} and \cref{eq:final-phelper-leq} follow by \cref{lem:nd-opt-contain}'s \cref{item:ND-contain-max-2} with $X\defeq C$, $X'\defeq Z$. Similarly, \cref{eq:nd-restrict-opt-lin-2} follows by \cref{lem:nd-opt-contain}'s \cref{item:ND-contain-max-2} with $X\defeq A$, $X'\defeq \ND{A}$. \Cref{eq:fewer-requirements} follows by applying the first inequality of \cref{lem:inclusion-opt} with $X\defeq \ND{A}, Y\defeq Z, Y'\defeq Z\setminus (B\setminus B')$. \Cref{eq:phi-helper-nd-opt} follows by applying \cref{lem:helper-perm} to \cref{eq:nd-restrict-opt-lin} with permutation $\phi$. 

\Cref{eq:phi-assumptions} follows by our assumptions on $\phi$. \Cref{eq:reunion} follows because by applying the second inequality of \cref{lem:inclusion-opt} with $X\defeq B', Y\defeq \ND{C}, Y'\defeq \ND{C}\setminus (B\setminus B')$.

Suppose $B'=B$. Then $B\setminus B'=\emptyset$, and so \cref{eq:fewer-requirements} and \cref{eq:reunion} are trivially equalities. Then \cref{eq:sim-nd-func} is an equality.

Suppose $\exists b<c: (b,c)^{\abs{\St}}\subseteq \supp[\Dany]$; note that $(b,c)^{\abs{\St}}\subseteq\supp[\phi\cdot \Dany]$, since such support must be invariant to permutation. Further suppose that $B'\subseteq C$ and that $\ND{C}\cap \prn{B\setminus B'}$ is non-empty. Then letting $X\defeq B', Y\defeq Z, Y'\defeq Z\setminus (B\setminus B')$ and noting that $\ND{\ND{Z}}=\ND{Z}$, apply \cref{lem:inclusion-opt} to \cref{eq:reunion} to conclude that \cref{eq:sim-nd-func} is strict.
\end{proof}

\begin{restatable}[Optimality probability superiority lemma]{lem}{optProbSup}\label{lem:opt-prob-superior}
Let $A,B,C\subsetneq \rewardVS$ be finite, and let $Z$ satisfy $\ND{C}\subseteq Z \subseteq C$. If $B$ contains a copy $B'$ of $\ND{A}$ via $\phi$ such that $\phi\cdot\prn{Z\setminus \prn{B\setminus B'}}=Z\setminus \prn{B\setminus B'}$, then $\phelper{A\geq C}[\Dany] \leqMost[][\DSetAny] \phelper{B \geq C}[\Dany]$.

If $B'\subseteq C$ and $\ND{C}\cap \prn{B\setminus B'}$ is non-empty, then the inequality is strict for all $\Diid\in\DSetBCiid$ and  $\phelper{A\geq C}[\Dany] \not\geqMost[][\DSetAny] \phelper{B \geq C}[\Dany]$.
\end{restatable}
\begin{proof}
Suppose $\Dany$ is such that $\phelper{B \geq C}[\Dany] < \phelper{A\geq C}[\Dany]$. 
\begin{align}
    \phelper{A\geq C}[\phi\cdot \Dany]&=\phelper{A\geq C}[\phi\inv\cdot \Dany]\label{eq:involution-1}\\
    &\leq \phelper{B\geq C}[\Dany]\label{eq:sim-apply-opt}\\
    &< \phelper{A\geq C}[\Dany]\label{eq:assumption-apply-leq}\\
    &\leq \phelper{B\geq C}[\phi\cdot \Dany].\label{eq:sim-apply-opt-2}
\end{align}

\Cref{eq:involution-1} holds because $\phi$ is an involution. \Cref{eq:sim-apply-opt} and \cref{eq:sim-apply-opt-2} hold by applying \cref{lem:sim-lin-func-opt} with permutation $\phi$. \Cref{eq:assumption-apply-leq} holds by assumption. Therefore, $\phelper{A\geq C}[\Dany] \leqMost[][\DSetAny] \phelper{B \geq C}[\Dany]$ by \cref{lem:helper-geq-most}.

Suppose $B'\subseteq C$ and $\ND{C}\cap \prn{B\setminus B'}$ is non-empty, and let $\Diid$ be any continuous distribution which distributes reward independently and identically across states. Let $\phi'\in\mdpPermGroup$.
\begin{align}
\phelper{A\geq C}[\phi'\cdot \Diid]&=\phelper{A\geq C}[\Diid]\label{eq:same-iid-dist-prob}\\
&<\phelper{B\geq C}[\phi\cdot \Diid]\label{eq:sim-iid-prob}\\
&=\phelper{A\geq C}[\phi'\cdot \Diid].\label{eq:same-iid-dist-prob-2}
\end{align}

\Cref{eq:same-iid-dist-prob} and \cref{eq:same-iid-dist-prob-2} hold because $\Diid$ distributes reward identically across states, $\forall \phi_x\in\mdpPermGroup: \phi_x\cdot \Diid=\Diid$. By \cref{lem:iid-contains}, $\exists b<c: (b,c)^{\abs{\St}}\subseteq \supp[\Diid]$. Therefore, apply \cref{lem:sim-lin-func-opt} to conclude that \cref{eq:sim-iid-prob} holds. 

Therefore, $\forall \phi'\in\mdpPermGroup:\phelper{A\geq C}[\phi'\cdot \Diid]<\phelper{B\geq C}[\phi'\cdot \Diid]$. In particular, $\phelper{A\geq C}[\Dany] \not\geqMost[][\DSetAny] \phelper{B \geq C}[\Dany]$ by \cref{def:ineq-most-dists}.
\end{proof}

\begin{restatable}[Limit probability inequalities which hold for most distributions]{lem}{limProbMost}\label{lem:lim-prob-most}
Let $I\subseteq \reals$, let $\distSet \subseteq \Delta(\rewardVS)$ be closed under permutation, and let $F_A,F_B,F_C$ be finite sets of vector functions $I\mapsto \rewardVS$. Let $\gamma$ be a limit point of $I$ such that $f_1(\D)\defeq \lim_{\gamma^*\to \gamma}\phelper{F_B(\gamma^*)\geq F_C(\gamma^*)}[\D],f_2(\D)\defeq \lim_{\gamma^*\to \gamma}\phelper{F_A(\gamma^*)\geq F_C(\gamma^*)}[\D]$ are well-defined for all $\D\in\distSet$.

Let $F_Z$ satisfy $\ND{F_C}\subseteq F_Z \subseteq F_C$. Suppose $F_B$ contains a copy of $F_A$ via $\phi$ such that $\phi\cdot \prn{F_Z\setminus \prn{F_B\setminus \phi\cdot F_A}}=F_Z\setminus \prn{F_B\setminus \phi\cdot F_A}$. Then $f_2(\distSet) \leqMost[][\distSet] f_1(\distSet)$.
\end{restatable}
\begin{proof}
Suppose $\D\in\distSet$ is such that $f_2(\D)>f_1(\D)$. 
\begin{align}
    f_2\prn{\phi\cdot \D}&=f_2\prn{\phi\inv\cdot \D}\label{eq:involute-f1}\\
    &\defeq \lim_{\gamma^*\to \gamma}\phelper{F_A(\gamma^*)\geq F_C(\gamma^*)}[\phi\inv\cdot \D]\\
    &\leq \lim_{\gamma^*\to \gamma}\phelper{F_B(\gamma^*)\geq F_C(\gamma^*)}[\D]\label{eq:FB-geq}\\
    &<\lim_{\gamma^*\to \gamma}\phelper{F_A(\gamma^*)\geq F_C(\gamma^*)}[\D]\label{eq:assumed-ineq}\\
    &\leq  \lim_{\gamma^*\to \gamma}\phelper{F_B(\gamma^*)\geq F_C(\gamma^*)}[\phi\cdot \D]\label{eq:FB-geq-2}\\
    &\eqdef f_1\prn{\phi\cdot \D}.
\end{align}
By the assumption that $\distSet$ is closed under permutation and $f_2$ is well-defined for all $\D\in\distSet$, $f_2(\phi\cdot \D)$ is well-defined. \Cref{eq:involute-f1} follows since $\phi=\phi\inv$ because $\phi$ is an involution. For all $\gamma^*\in I$, let $A\defeq F_A(\gamma^*), B\defeq F_B(\gamma^*), C\defeq F_C(\gamma^*), Z\defeq F_Z(\gamma^*)$ (by \cref{def:nd-vec-func}, $\ND{C}\subseteq Z \subseteq C$). Since $\phi\cdot A\subseteq B$ by assumption, and since $\ND{A}\subseteq A$, $B$ also contains a copy of $\ND{A}$ via $\phi$. Furthermore, $\phi\cdot \prn{Z\setminus \prn{B\setminus \phi\cdot A}}=Z\setminus \prn{B\setminus \phi\cdot A}$ (by assumption), and so apply  \cref{lem:sim-lin-func-opt} to conclude that $\phelper{F_A(\gamma^*)\geq F_C(\gamma^*)}[\phi\inv\cdot \D]\leq \phelper{F_B(\gamma^*)\geq F_C(\gamma^*)}[\D]$. Therefore, the limit inequality \cref{eq:FB-geq} holds. \Cref{eq:assumed-ineq} follows because we assumed that $f_1(\D)<f_2(\D)$. \Cref{eq:FB-geq-2} holds by reasoning similar to that given for \cref{eq:FB-geq}.

Therefore, $f_2(\D)>f_1(\D)$ implies that $f_2\prn{\phi\cdot \D}<f_1\prn{\phi\cdot \D}$, and so apply \cref{lem:helper-geq-most} to conclude that $f_2(\D) \leqMost[][\distSet] f_1(\D)$.
\end{proof}

\subsubsection{\texorpdfstring{$\Fnd$}{Non-dominated visit distribution} results} 
  
\transferDiscount
\begin{proof}
Let $R$ be any reward function. Suppose $\gamma^*\in(0,1)$ and construct $R'(s)\defeq \OptVf{s, \gamma}- \gamma^* \max_{a\in\A}\E{s'\sim T(s,a)}{\OptVf{s',\gamma}}$. 

Let $\pi\in\Pi$ be any policy. By the definition of optimal policies, $\pi\in\optPi[R',\gamma^*]$ iff for all $s$:
\begin{align}
    R'(s) + \gamma^* \E{s' \sim T\prn{s,\pi(s)}}{\OptVf[R']{s',\gamma^*}}&=R'(s) + \gamma^* \max_{a\in\A}\E{s' \sim T\prn{s,a}}{\OptVf[R']{s',\gamma^*}}\\
    R'(s) + \gamma^* \E{s' \sim T\prn{s,\pi(s)}}{\OptVf{s',\gamma}}&=R'(s) + \gamma^* \max_{a\in\A}\E{s' \sim T\prn{s,a}}{\OptVf{s',\gamma}}\label{eq:opt-val-fn-equality}\\
    \gamma^*\E{s' \sim T\prn{s,\pi(s)}}{\OptVf{s',\gamma}}&=\gamma^*\max_{a\in\A}\E{s'\sim T(s,a)}{\OptVf{s',\gamma}}\label{eq:subst-r'}\\
    \E{s' \sim T\prn{s,\pi(s)}}{\OptVf{s',\gamma}}&=\max_{a\in\A}\E{s'\sim T(s,a)}{\OptVf{s',\gamma}}.\label{eq:div-gamma*}
\end{align} 

By the Bellman equations, $R'(s)=\OptVf[R']{s, \gamma^*}- \gamma^* \max_{a\in\A}\E{s'\sim T(s,a)}{\OptVf[R']{s',\gamma^*}}$. By the definition of $R'$, $\OptVf[R']{\cdot,\gamma^*}=\OptVf[R]{\cdot, \gamma}$ must be the unique solution to the Bellman equations for $R'$ at $\gamma^*$. Therefore, \cref{eq:opt-val-fn-equality} holds. \Cref{eq:subst-r'} follows by plugging in $R'\defeq \OptVf{s, \gamma}- \gamma^* \max_{a\in\A}\E{s'\sim T(s,a)}{\OptVf{s',\gamma}}$ to \cref{eq:opt-val-fn-equality} and doing algebraic manipulation. \Cref{eq:div-gamma*} follows because $\gamma^*>0$.  

\Cref{eq:div-gamma*} shows that $\pi\in\optPi[R',\gamma^*]$ iff $\forall s: \E{s'\sim T(s,\pi(s))}{\OptVf{s',\gamma}} = \max_{a\in\A} \E{s'\sim T(s,a)}{\OptVf{s',\gamma}}$. That is, $\pi\in\optPi[R',\gamma^*]$ iff $\pi\in\optPi[R,\gamma]$.
\end{proof} 

\begin{restatable}[Evaluating sets of visit distribution functions at $\gamma$]{definition}{evalFDisc}\label{def:eval-f-discount}
For $\gamma \in (0,1)$, define $\F(s,\gamma)\defeq \set{\f(\gamma) \mid \f \in \F(s)}$ and $\Fnd(s,\gamma)\defeq \set{\f(\gamma) \mid \f \in \Fnd(s)}$. If $F\subseteq \F(s)$, then $F(\gamma) \defeq \set{\f(\gamma) \mid \f \in F}$.
\end{restatable}

\begin{restatable}[Non-domination across $\gamma$ values for expectations of visit distributions]{lem}{ndSetFGamma}\label{lem:nd-gamma-F-subset}
Let $\Delta_d \in \Delta\prn{\rewardVS}$ be any state distribution and let $F\defeq\set{\E{s_d\sim\Delta_d}{\fpi{s_d}} \mid \pi \in \Pi}$. $\f \in \ND{F}$ iff $\forall \gamma^*\in(0,1): \f(\gamma^*)\in \ND{F(\gamma^*)}$.
\end{restatable}
\begin{proof}
Let $\fpi{}\in\ND{F}$ be strictly optimal for reward function $R$ at discount rate $\gamma\in(0,1)$:
\begin{align}
    \fpi{}(\gamma)^\top \rf > \max_{\fpi[\pi']{}\in F\setminus \set{\fpi{}}} \fpi[\pi']{}(\gamma)^\top \rf.
\end{align}

Let $\gamma^*\in(0,1)$. By \cref{transferDiscount}, we can produce $R'$ such that $\optPi[R',\gamma^*]=\optPi$. Since the optimal policy sets are equal, \cref{lem:opt-pol-visit-iff} implies that
\begin{align}
    \fpi{}(\gamma^*)^\top \rf' > \max_{\fpi[\pi']{}\in F\setminus \set{\fpi{}}} \fpi[\pi']{}(\gamma^*)^\top \rf'.
\end{align}

Therefore, $\fpi{}(\gamma^*)\in\ND{F(\gamma^*)}$. 

The reverse direction follows by the definition of $\ND{F}$.
\end{proof}

\begin{restatable}[$\forall \gamma \in (0,1): \dbf\in \Fnd(s,\gamma)$ iff $\dbf\in\ND{\F(s,\gamma)}$]{lem}{invariantFndDiscount}\label{lem:nd-relation} 
\end{restatable} 
\begin{proof}
By \cref{def:eval-f-discount}, $\Fnd(s,\gamma)\defeq \set{\f(\gamma) \mid \f \in \ND{\F(s)}}$. By applying \cref{lem:nd-gamma-F-subset} with $\Delta_d\defeq \unitvec$, $\f \in \ND{\F(s)}$ iff $\forall \gamma\in(0,1): \f(\gamma)\in \ND{\F(s,\gamma)}$. 
\end{proof}

\optVfFndRestrict
\begin{proof}
$\ND{\F(s,\gamma)}=\Fnd(s,\gamma)$ by \cref{lem:nd-relation}, so apply \cref{cor:nd-func-indif} with $X \defeq \F(s,\gamma)$. 
\end{proof} 

\subsection{Some actions have greater probability of being optimal}
\begin{restatable}[Optimal policy shift bound]{lem}{optPiBound}\label{lem:opt-pol-shift-bound}
For fixed $R$, $\optPi$ can take on at most $(2\abs{\St}+1)\sum_s \binom{\abs{\F(s)}}{2}$ distinct values over $\gamma \in (0,1)$.
\end{restatable} 
\begin{proof}
By \cref{lem:opt-pol-visit-iff}, $\optPi$ changes value iff there is a change in optimality status for some visit distribution function at some state. \citet{lippman1968set} showed that two visit distribution functions can trade off optimality status at most $2\abs{\St}+1$ times. At each state $s$, there are $\binom{\abs{\F(s)}}{2}$ such pairs. 
\end{proof} 

\begin{restatable}[Optimality probability's limits exist]{prop}{muConverge} \label{prop:opt-prob-converge} Let $F\subseteq \F(s)$. $\optprob[\Dany]{F,0}=\lim_{\gamma\to 0} \optprob[\Dany]{F,\gamma}$ and $\optprob[\Dany]{F,1}=\lim_{\gamma\to 1} \optprob[\Dany]{F,\gamma}$. 
\end{restatable}
\begin{proof}
First consider the limit as $\gamma \to 1$. Let $\Dany$ have probability measure $F_\text{any}$, and define $\delta(\gamma)\defeq F_\text{any}\prn{\set{R\in\rewardSpace \mid \exists \gamma^* \in [\gamma,1): \optPi[R,\gamma^*]\neq \optPi[R,1]}}$. Since $F_\text{any}$ is a probability measure, $\delta(\gamma)$ is bounded $[0,1]$, and $\delta(\gamma)$ is monotone decreasing. Therefore, $\lim_{\gamma \to 1} \delta(\gamma)$ exists. 

If $\lim_{\gamma \to 1} \delta(\gamma)>0$, then there exist reward functions whose optimal policy sets $\optPi$ never converge (in the discrete topology on sets) to $\optPi[R,1]$, contradicting \cref{lem:opt-pol-shift-bound}. So $\lim_{\gamma \to 1} \delta(\gamma)=0$.

By the definition of optimality probability (\cref{def:prob-opt}) and of $\delta(\gamma)$,  $|\optprob[\Dany]{F,\gamma}-\optprob[\Dany]{F,1}|\leq \delta(\gamma)$. Since $\lim_{\gamma\to 1} \delta(\gamma)=0$, $\lim_{\gamma\to 1}\optprob[\Dany]{F,\gamma}=\optprob[\Dany]{F,1}$.

A similar proof shows that $\lim_{\gamma\to 0} \optprob[\Dany]{F,\gamma}=\optprob[\Dany]{F,0}$.
\end{proof} 

\begin{restatable}[Optimality probability identity]{lem}{optProbID}\label{lem:opt-prob-id}
Let $\gamma\in(0,1)$ and let $F\subseteq \F(s)$. 
\begin{equation}
\optprob[\Dany]{F,\gamma}=\phelper{F(\gamma)\geq \F(s,\gamma)}=\phelper{F(\gamma)\geq \Fnd(s,\gamma)}.
\end{equation}
\end{restatable}
\begin{proof}
Let $\gamma\in(0,1)$.
\begin{align}
    \optprob[\Dany]{F,\gamma}&\defeq \prob[R\sim\Dany]{\exists \f^\pi \in F: \pi\in\optPi}\\
    &=\E{\rf \sim\Dany}{ \indic{\max_{\f\in F} \f(\gamma)^\top \rf = \max_{\f'\in\F(s)} \f'(\gamma)^\top \rf}}\label{eq:F-probability}\\
    &=\E{\rf \sim\Dany}{ \indic{\max_{\f\in F} \f(\gamma)^\top \rf = \max_{\f'\in\Fnd(s)} \f'(\gamma)^\top \rf}}\label{eq:Fnd-probability}\\
    &\eqdef \phelper{F(\gamma)\geq \Fnd(s,\gamma)}.
\end{align}
\Cref{eq:F-probability} follows because \cref{lem:opt-pol-visit-iff} shows that $\pi$ is optimal iff it induces an optimal visit distribution  $\f$ at every state. \Cref{eq:Fnd-probability} follows because $\forall \rf \in \rewardVS: \max_{\f'\in\F(s)} \f'(\gamma)^\top \rf= \max_{\f'\in\Fnd(s)} \f'(\gamma)^\top \rf$ by \cref{cor:opt-vf-restrict-fnd}.
\end{proof}

\subsection{Basic properties of \texorpdfstring{$\pwrNoDist$}{power}}

\begin{restatable}[$\pwrNoDist$ identities]{lem}{lemPowEQ}\label{lem:power-id}
Let $\gamma \in (0,1)$.
\begin{align}
    \pwr[s,\gamma]&=\E{\rf\sim\Dbd}{\max_{\f\in \Fnd(s)} \frac{1-\gamma}{\gamma}\prn{\f(\gamma)-\unitvec}^\top \rf}\label{eq:pwr-def-f}\\
    &= \dfrac{1-\gamma}{\gamma}\E{\rf\sim\Dbd}{\OptVf{s,\gamma}-R(s)}\\
    &=  \dfrac{1-\gamma}{\gamma}\prn{\vavg-\E{R\sim \Dbd}{R(s)}} \\
    &=\E{R\sim\Dbd}{\max_{\pi\in\Pi} \E{s'\sim T\prn{s,\pi(s)}}{\prn{1-\gamma}V^\pi_R\prn{s',\gamma}}}.\label{eq:pwr-def-avg-discounted}
\end{align}
\end{restatable}
\begin{proof}
\begin{align}
    \pwrNoDist_{\Dbd}(s,\gamma)&\defeq\E{\rf\sim\Dbd}{\max_{\f\in \F(s)} \frac{1-\gamma}{\gamma}\prn{\f(\gamma)-\unitvec}^\top \rf}\\
    &= \E{\rf\sim\Dbd}{\max_{\f\in \Fnd(s)} \frac{1-\gamma}{\gamma}\prn{\f(\gamma)-\unitvec}^\top \rf}\label{eq:pwr-fnd-restrict}\\
    &=\E{\rf\sim\Dbd}{\max_{\f\in \F(s)} \frac{1-\gamma}{\gamma}\prn{\f(\gamma)-\unitvec}^\top \rf}\\
    &= \dfrac{1-\gamma}{\gamma}\E{\rf\sim\Dbd}{\OptVf{s,\gamma}-R(s)}\label{eq:pwr-vf-convert}\\
    &=\dfrac{1-\gamma}{\gamma}\prn{\vavg-\E{R\sim \Dbd}{R(s)}}\label{eq:vavg-def-pwr-def}\\
    &= \E{\rf\sim\Dbd}{\max_{\pi\in\Pi} \E{s'\sim T\prn{s,\pi(s)}}{\prn{1-\gamma}\fpi[\pi]{s'}(\gamma)^\top\rf}}\label{eq:f-expand-pwr-def}\\
    &= \E{R\sim\Dbd}{\max_{\pi\in\Pi} \E{s'\sim T\prn{s,\pi(s)}}{\prn{1-\gamma}V^\pi_R\prn{s',\gamma}}}.\label{eq:pwr-avg-normalized}
\end{align}

\Cref{eq:pwr-fnd-restrict} follows from \cref{cor:opt-vf-restrict-fnd}. \Cref{eq:pwr-vf-convert} follows from the dual formulation of optimal value functions. \Cref{eq:vavg-def-pwr-def} holds by the definition of $\vavg$ (\cref{def:vavg}). \Cref{eq:f-expand-pwr-def} holds because $\fpi{s}(\gamma) = \unitvec +\gamma \E{s'\sim T\prn{s,\pi(s)}}{\fpi[\pi]{s'}(\gamma)}$ by the definition of a visit distribution function (\cref{def:visit}).
\end{proof}

\begin{restatable}[Discount-normalized value function]{definition}{discValue}\label{def:disc-value}
Let $\pi$ be a policy, $R$ a reward function, and $s$ a state. For $\gamma \in [0,1]$, $\VfNorm[\pi]{s,\gamma}\defeq \lim_{\gamma^*\to \gamma}(1-\gamma^*) V^\pi_R(s,\gamma^*)$.
\end{restatable}

\begin{restatable}[Normalized value functions have uniformly bounded derivative]{lem}{normValueLip}\label{lem:norm-value-lip}
There exists $K\geq0$ such that for all reward functions $\rf\in \rewardVS$, $\sup_{\substack{s\in\St, \pi\in\Pi,\gamma\in [0,1]}}\abs{\frac{d}{d\gamma} \VfNorm[\pi]{s,\gamma}} \leq K\lone{\rf}$.
\end{restatable}
\begin{proof}
Let $\pi$ be any policy, $s$ a state, and $R$ a reward function. Since $\VfNorm[\pi]{s,\gamma}=\lim_{\gamma^*\to\gamma}(1-\gamma^*)\fpi{s}(\gamma^*)^\top \rf$, $\frac{d}{d\gamma}\VfNorm[\pi]{s,\gamma}$ is controlled by the behavior of $\lim_{\gamma^*\to\gamma}(1-\gamma^*)\fpi{s}(\gamma^*)$. We show that this function's gradient is bounded in infinity norm.

By \cref{f-rat}, $\fpi{s}(\gamma)$ is a multivariate rational function on $\gamma$. Therefore, for any state $s'$,  $\fpi{s}(\gamma)^\top\unitvec[s']=\frac{P(\gamma)}{Q(\gamma)}$ in reduced form. By \cref{prop:visit-dist-prop}, $0\leq \fpi{s}(\gamma)^\top\unitvec[s']\leq \frac{1}{1-\gamma}$. Thus, $Q$ may only have a root of multiplicity 1 at $\gamma=1$, and $Q(\gamma)\neq 0$ for $\gamma\in[0,1)$. Let $f_{s'}(\gamma)\defeq (1-\gamma)\fpi{s}(\gamma)^\top\unitvec[s']$.

If $Q(1)\neq 0$, then the derivative $f_{s'}'(\gamma)$ is bounded on $\gamma\in[0,1)$ because the polynomial $(1-\gamma)P(\gamma)$ cannot diverge on a bounded domain. 

If $Q(1)=0$, then factor out the root as $Q(\gamma)=(1-\gamma)Q^*(\gamma)$. 
\begin{align}
    f_{s'}'(\gamma)&=\frac{d}{d\gamma}\prn{\frac{(1-\gamma)P(\gamma)}{Q(\gamma)}}\\
    &=\frac{d}{d\gamma}\prn{\frac{P(\gamma)}{Q^*(\gamma)}}\\
    &=\frac{P'(\gamma)Q^*(\gamma)-(Q^*)'(\gamma)P(\gamma)}{(Q^*(\gamma))^2}.
\end{align}

Since $Q^*(\gamma)$ is a polynomial with no roots on $\gamma\in[0,1]$, $f_{s'}'(\gamma)$ is bounded on $\gamma\in[0,1)$.

Therefore, whether or not $Q(\gamma)$ has a root at $\gamma=1$, $f_{s'}'(\gamma)$ is bounded on $\gamma\in[0,1)$. Furthermore, $\sup_{\gamma\in[0,1)}\linfty{\nabla(1-\gamma)\fpi{s}(\gamma)}=\sup_{\gamma\in[0,1)}\max_{s'\in\St} \abs{f_{s'}'(\gamma)}$ is finite since there are only finitely many states.

There are finitely many $\pi\in\Pi$, and finitely many states $s$, and so there exists some $K'$ such that $\sup_{\substack{s\in\St,\\ \pi\in\Pi,\gamma\in [0,1)}} \linfty{\nabla(1-\gamma)\fpi{s}(\gamma)}\leq K'$. Then $\lone{\nabla(1-\gamma)\fpi{s}(\gamma)}\leq \abs{\St}K'\eqdef K$. 
\begin{align}
    \sup_{\substack{s\in\St,\\ \pi\in\Pi,\gamma\in [0,1)}} \abs{\frac{d}{d\gamma}V^\pi_{R,\text{norm}}\prn{s,\gamma}}\defeq\,&\sup_{\substack{s\in\St,\\ \pi\in\Pi,\gamma\in [0,1)}} \abs{\frac{d}{d\gamma}\lim_{\gamma^*\to\gamma}(1-\gamma^*)V^\pi_{R}\prn{s,\gamma^*}}\\
    =\,&\sup_{\substack{s\in\St,\\ \pi\in\Pi,\gamma\in [0,1)}} \abs{\frac{d}{d\gamma}(1-\gamma)V^\pi_{R}\prn{s,\gamma}}\label{eq:lim-disappear}\\
    =\,&\sup_{\substack{s\in\St,\\ \pi\in\Pi,\gamma\in [0,1)}} \abs{\nabla (1-\gamma)\fpi{s}(\gamma)^\top\rf}\label{eq:on-pol-rat}\\
    \leq\, & \sup_{\substack{s\in\St,\\ \pi\in\Pi,\gamma\in [0,1)}} \lone{\nabla (1-\gamma)\fpi{s}(\gamma)}\lone{\rf}\label{eq:cs}\\
    \leq\, & K\lone{\rf}.\label{eq:bounded}
\end{align}

\Cref{eq:lim-disappear} holds because $\Vf[\pi]{s,\gamma}$ is continuous on $\gamma\in[0,1)$ by \cref{smoothOnPol}. \Cref{eq:cs} holds by the Cauchy-Schwarz inequality.

Since $\abs{\frac{d}{d\gamma}V^\pi_{R,\text{norm}}\prn{s,\gamma}}$ is bounded for all $\gamma \in [0,1)$, \cref{eq:bounded} also holds for $\gamma\to 1$. 
\end{proof}

\ContPower*
\begin{proof}
Let $b,c$ be such that $\supp[\Dbd]\subseteq [b,c]^{\abs{\St}}$. For any $\rf\in\supp[\Dbd]$ and $\pi\in\Pi$, $\VfNorm[\pi]{s,\gamma}$ has Lipschitz constant $K\lone{\rf}\leq K \abs{\St}\linfty{\rf}\leq K\abs{\St}\max(\abs{c},\abs{b})$ on $\gamma\in(0,1)$ by \cref{lem:norm-value-lip}.

For $\gamma\in(0,1)$, $\pwr[s,\gamma] = \E{R\sim\Dbd}{\max_{\pi\in\Pi} \E{s'\sim T\prn{s,\pi(s)}}{(1-\gamma)V^\pi_R\prn{s',\gamma}}}$ by \cref{eq:pwr-avg-normalized}. The expectation of the maximum of a set of functions which share a Lipschitz constant, also shares the Lipschitz constant. This shows that $\pwr[s,\gamma]$ is Lipschitz continuous on $\gamma\in(0,1)$. Thus, its limits are well-defined as $\gamma\to 0$ and $\gamma\to 1$. So it is Lipschitz continuous on the closed unit interval.
\end{proof}

\maxPwrGeneral*
\begin{proof}
Let $\gamma \in (0,1)$.
\begin{align}
    \pwr[s,\gamma][\Dbd]&= \E{R\sim \Dbd}{\max_{\pi\in\Pi} \E{s'\sim T(s,\pi(s))}{(1-\gamma)\OptVf{s',\gamma}}}\label{eq:pwr-id}\\
    &\leq \E{R\sim \Dbd}{\max_{\pi\in\Pi} \E{s'\sim T(s,\pi(s))}{(1-\gamma)\geom[\max_{s''\in\St}R(s'')]}}\label{eq:max-vfn}\\
    &= \E{R\sim \Dbd}{\max_{s''\in\St}R(s'')}.
\end{align}

\Cref{eq:pwr-id} follows from \cref{lem:power-id}. \Cref{eq:max-vfn} follows because $\OptVf{s',\gamma}\leq \geom[\max_{s''\in\St}R(s'')]$, as no policy can do better than achieving maximal reward at each time step. Taking limits, the inequality holds for all $\gamma\in[0,1]$.

Suppose that $s$ can deterministically reach all states in one step and all states are {\stateEnd}s. Then \cref{eq:max-vfn} is an equality for all $\gamma\in(0,1)$, since for each $R$, the agent can select an action which deterministically transitions to a state with maximal reward. Thus the equality holds for all $\gamma\in[0,1]$.
\end{proof}

\begin{restatable}[Lower bound on current $\pwrNoDist$ based on future $\pwrNoDist$]{lem}{FuturePower}\label{lem:future-power}
\begin{align}
\pwr[s,\gamma]\geq (1-\gamma)\min_a \E{\substack{s'\sim T(s,a),\\R\sim\Dbd}}{R(s')} + \gamma\max_a \E{s'\sim T(s,a)}{\pwr[s',\gamma]}.
\end{align}
\end{restatable}
\begin{proof} 
Let $\gamma \in (0,1)$ and let $a^*\in \argmax_a \E{s'\sim T\prn{s,a}}{\pwr[s',\gamma]}$.
\begin{align}
    &\pwr[s,\gamma]\\
    =\,&(1-\gamma)\E{R\sim\Dbd}{\max_a \E{s'\sim T\prn{s,a}}{\OptVf{s',\gamma}}}\label{eq:pwr-id-app}\\
    \geq\,& (1-\gamma) \max_a \E{s'\sim T\prn{s,a}}{\E{R\sim\Dbd}{\OptVf{s',\gamma}}}\label{eq:max-expect-ineq}\\
    =\,&(1-\gamma)\max_a \E{s'\sim T\prn{s,a}}{\vavg[s',\gamma]}\\
    =\,&(1-\gamma) \max_a \E{s'\sim T\prn{s,a}}{\E{R\sim \Dbd}{R(s')}+\frac{\gamma}{1-\gamma}\pwr[s',\gamma]}\label{eq:vavg-identity}\\
    \geq\,&(1-\gamma) \E{s'\sim T\prn{s,a^*}}{\E{R\sim \Dbd}{R(s')}+\frac{\gamma}{1-\gamma}\pwr[s',\gamma]}\\
    \geq\,&(1-\gamma)\min_a \E{\substack{s'\sim T(s,a),\\R\sim\Dbd}}{R(s')} + \gamma \E{s'\sim T\prn{s,a^*}}{\pwr[s',\gamma]}.\label{eq:min-reward-max-power}
\end{align}

\Cref{eq:pwr-id-app} holds by \cref{lem:power-id}. \Cref{eq:max-expect-ineq} follows because $\E{x\sim X}{\max_a f(a,x)}\geq \max_a \E{x\sim X}{f(a,x)}$ by Jensen's inequality, and \cref{eq:vavg-identity} follows by \cref{lem:power-id}. 

The inequality also holds when we take the limits  $\gamma \to 0$ or $\gamma \to 1$.
\end{proof}

\smooth*
\begin{proof}
Suppose $\gamma\in[0,1]$. First consider the case where $\pwr[s,\gamma]\geq \pwr[s',\gamma]$.
\begin{align}
\pwr[s',\gamma]&\geq (1-\gamma)\min_a \E{\substack{s_x\sim T(s',a),\\R\sim\Dbd}}{R(s_x)} + \gamma\max_a \E{s_x\sim T(s',a)}{\pwr[s_x,\gamma]}\label{eq:lb-pwr-next}\\
&\geq (1-\gamma)b + \gamma \pwr[s,\gamma].\label{eq:lb-pwr-next-min}
\end{align}

\Cref{eq:lb-pwr-next} follows by \cref{lem:future-power}. \Cref{eq:lb-pwr-next-min} follows because reward is lower-bounded by $b$ and because  $s'$ can reach $s$ in one step with probability 1. 
\begin{align}
    \abs{\pwr[s,\gamma]-\pwr[s',\gamma]}&=\pwr[s,\gamma]-\pwr[s',\gamma]\label{eq:no-abs}\\
    &\leq \pwr[s,\gamma]-\prn{(1-\gamma)b + \gamma \pwr[s,\gamma]}\label{eq:lb-pwr-abs-diff}\\
    &=(1-\gamma)\prn{\pwr[s,\gamma] - b}\\
    &\leq(1-\gamma)\prn{\E{R\sim \Dbd}{\max_{s''\in\St}R(s'')} - b}\label{eq:ub-pwr-s}\\
    &\leq (1-\gamma)(c-b).\label{eq:ub-pwr-abs-diff}
\end{align}

\Cref{eq:no-abs} follows because $\pwr[s,\gamma]\geq \pwr[s',\gamma]$. \Cref{eq:lb-pwr-abs-diff} follows by \cref{eq:lb-pwr-next-min}. \Cref{eq:ub-pwr-s} follows by \cref{lem:max-power-general}. \Cref{eq:ub-pwr-abs-diff} follows because reward under $\Dbd$ is upper-bounded by $c$.

The case where $\pwr[s,\gamma]\leq \pwr[s',\gamma]$ is similar, leveraging the fact that $s$ can also reach $s'$ in one step with probability 1.
\end{proof}

\subsection{Seeking \texorpdfstring{$\pwrNoDist$}{power} is often more probable under optimality}

\subsubsection{Keeping options open tends to be \texorpdfstring{$\pwrNoDist$}{POWER}-seeking and tends to be optimal}

\begin{restatable}[Normalized visit distribution function]{definition}{normVisitFn}\label{def:norm-visit-fn}
Let $\f:[0,1)\to \rewardVS$ be a vector function. For $\gamma\in[0,1]$, $\NormF{\f,\gamma}\defeq \lim_{\gamma^*\to\gamma} (1-\gamma^*)\f(\gamma^*)$ (this limit need not exist for arbitrary $\f$). If $F$ is a set of such $\f$, then $\NormF{F,\gamma}\defeq \set{\NormF{\f,\gamma}\mid \f \in F}$.
\end{restatable}

\begin{remark}
$\RSD=\NormF{\F(s),1}$.
\end{remark}

\begin{restatable}[Normalized visit distribution functions are continuous]{lem}{contNormFMix}\label{lem:cont-norm-f}
Let $\Delta_s\in \Delta(\St)$ be a state probability distribution, let $\pi\in\Pi$, and let $\f^*\defeq \E{s\sim \Delta_s}{\fpi{s}}$. $\NormF{\f^*, \gamma}$ is continuous on $\gamma\in [0,1]$.
\end{restatable}
\begin{proof}
\begin{align}
    \NormF{\f^*,\gamma}&\defeq \lim_{\gamma^*\to\gamma} (1-\gamma^*)\E{s\sim \Delta_s}{\fpi{s}(\gamma^*)}\\
    &=\E{s\sim \Delta_s}{\lim_{\gamma^*\to\gamma} (1-\gamma^*)\fpi{s}(\gamma^*)}\label{eq:expect-lim-swap}\\
    &\eqdef \E{s\sim \Delta_s}{\NormF{\fpi{s},\gamma}}.\label{eq:expect-cont}
\end{align}
\Cref{eq:expect-lim-swap} follows because the expectation is over a finite set. Each $\fpi{s}\in\F(s)$ is continuous on $\gamma\in[0,1)$ by \cref{f-rat}, and $\lim_{\gamma^*\to 1} (1-\gamma^*) \fpi{s}(\gamma^*)$ exists because {\rsd}s are well-defined \citep{puterman_markov_2014}. Therefore, each $\NormF{\fpi{s},\gamma}$ is continuous on $\gamma\in[0,1]$. Lastly, \cref{eq:expect-cont}'s expectation over finitely many continuous functions is itself continuous. 
\end{proof}

\begin{restatable}[Non-domination of normalized visit distribution functions]{lem}{ndNormalVisit}\label{lem:nd-norm-visit}
Let $\Delta_s\in \Delta(\St)$ be a state probability distribution and let $F\defeq \set{\E{s\sim \Delta_s}{\fpi{s}}\mid\pi\in\Pi}$. For all $\gamma\in[0,1]$, $\ND{\NormF{F,\gamma}}\subseteq \NormF{\ND{F},\gamma}$, with equality when $\gamma\in(0,1)$.
\end{restatable}
\begin{proof}
Suppose $\gamma\in (0,1)$.
\begin{align}
    \ND{\NormF{F,\gamma}}&=\ND{(1-\gamma)F(\gamma)}\label{eq:apply-cont-norm-1}\\
    &= (1-\gamma)\ND{F(\gamma)}\label{eq:nd-rescale-invariant}\\
    &= (1-\gamma)\prn{\ND{F}(\gamma)}\label{eq:nd-invariant-01}\\
    &= \NormF{\ND{F},\gamma}.\label{eq:apply-cont-norm-2}
\end{align}

\Cref{eq:apply-cont-norm-1} and \cref{eq:apply-cont-norm-2} follow by the continuity of $\NormF{\f,\gamma}$ (\cref{lem:cont-norm-f}). \Cref{eq:nd-rescale-invariant} follows by \cref{lem:pos-aff-nd-invar} \cref{item:invar-vectors}. \Cref{eq:nd-invariant-01} follows by \cref{lem:nd-gamma-F-subset}.

Let $\gamma=1$. Let $\dbf \in \ND{\NormF{F,1}}$ be strictly optimal for $\rf^*\in\rewardVS$. Then let $F_\dbf \subseteq F$ be the subset of $\f\in F$ such that $\NormF{\f,1}=\dbf$. 
\begin{align}
    \max_{\f\in F_\dbf} \NormF{\f,1}^\top \rf^* & > \max_{\f'\in F\setminus F_\dbf} \NormF{\f',1}^\top \rf^*.\label{eq:norm-f-ineq}
\end{align}

Since $\NormF{\f,1}$ is continuous at $\gamma=1$ (\cref{lem:cont-norm-f}), $\x^\top \rf^*$ is continuous on $\x\in\rewardVS$, and $F$ is finite, \cref{eq:norm-f-ineq} holds for some $\gamma^*\in(0,1)$ sufficiently close to $\gamma=1$. By \cref{lem:all-rf-max-nd}, at least one $\f\in F_\dbf$ is an element of $\ND{F(\gamma^*)}$. Then by \cref{lem:nd-gamma-F-subset}, $\f\in\ND{F}$. We conclude that $\ND{\NormF{F,1}}\subseteq \NormF{\ND{F},1}$. 

The case for $\gamma=0$ proceeds similarly.
\end{proof}

\begin{restatable}[$\pwrNoDist$ limit identity]{lem}{pwrLimit}\label{lem:pwr-limit}
Let $\gamma\in[0,1]$.
\begin{align}
    \pwr[s,\gamma]&=\E{\rf\sim\Dbd}{\max_{\f\in \Fnd(s)} \lim_{\gamma^*\to\gamma}\frac{1-\gamma^*}{\gamma^*}\prn{\f(\gamma^*)-\unitvec}^\top \rf}.
\end{align}
\end{restatable}
\begin{proof}
Let $\gamma\in[0,1]$.
\begin{align}
    \pwr[s,\gamma]&= \lim_{\gamma^*\to\gamma}\pwr[s,\gamma^*]\label{eq:pwr-cont-lim}\\
    &=\lim_{\gamma^*\to\gamma} \E{\rf\sim\Dbd}{\max_{\f\in \Fnd(s)} \frac{1-\gamma^*}{\gamma^*}\prn{\f(\gamma^*)-\unitvec}^\top \rf}\label{eq:pwr-lim-def}\\
    &=\E{\rf\sim\Dbd}{\lim_{\gamma^*\to\gamma}\max_{\f\in \Fnd(s)} \frac{1-\gamma^*}{\gamma^*}\prn{\f(\gamma^*)-\unitvec}^\top \rf}\label{eq:pwr-dom-conv}\\
    &=\E{\rf\sim\Dbd}{\max_{\f\in \Fnd(s)} \lim_{\gamma^*\to\gamma}\frac{1-\gamma^*}{\gamma^*}\prn{\f(\gamma^*)-\unitvec}^\top \rf}.\label{eq:max-cont-lim}
\end{align}
\Cref{eq:pwr-cont-lim} follows because $\pwr[s,\gamma]$ is continuous on $\gamma\in[0,1]$ by \cref{thm:cont-power}. \Cref{eq:pwr-lim-def} follows by \cref{lem:power-id}. 

For $\gamma^*\in(0,1)$, let $f_{\gamma^*}(\rf)\defeq \max_{\f\in \Fnd(s)} \frac{1-\gamma^*}{\gamma^*}\prn{\f(\gamma^*)-\unitvec}^\top \rf$. For any sequence $\gamma_n\to\gamma$, $\prn{f_{\gamma_n}}_{n=1}^\infty$ is a sequence of functions which are piecewise linear on $\rf \in \rewardVS$, which means they are continuous and therefore measurable. Since \cref{f-rat} shows that each $\f\in\Fnd(s)$ is multivariate rational on $\gamma^*$ (and therefore continuous on $\gamma^*$), $\set{f_{\gamma_n}}_{n=1}^\infty$ converges pointwise to limit function $f_\gamma$. Furthermore,  $\abs{\OptVf{s,\gamma_n}-R(s)}\leq \frac{\gamma}{1-\gamma_n}\linfty{R}$, and so $\abs{f_{\gamma_n}(\rf)}=\abs{\frac{1-\gamma_n}{\gamma_n}(\OptVf{s,\gamma_n}-R(s))}\leq g(\rf)\leq \linfty{\rf}\eqdef g(\rf)$, which is measurable. Therefore, apply Lebesgue's dominated convergence theorem to conclude that \cref{eq:pwr-dom-conv} holds. \Cref{eq:max-cont-lim} holds because $\max$ is a continuous function.
\end{proof}

\begin{restatable}[Lemma for $\pwrNoDist$ superiority]{lem}{morePowerPrefixmoreOptions}\label{lem:more-power-prefix}
Let $\Delta_1,\Delta_2\in \Delta\prn{\St}$ be state probability distributions. For $i=1,2$, let $F_{\Delta_i}\defeq \set{\gamma\inv\E{s_i\sim \Delta_i}{\fpi{s_i}-\unitvec[s_i]}\mid \pi\in\Pi}$. Suppose $F_{\Delta_2}$ contains a copy of $\ND{F_{\Delta_1}}$ via $\phi$. Then $\forall \gamma\in[0,1]:\E{s_1\sim \Delta_1}{\pwr[s_1,\gamma]}\leqMost[][\DSetBd] \E{s_2\sim \Delta_2}{\pwr[s_2,\gamma]}$. 

If $\ND{F_{\Delta_2}}\setminus \phi\cdot \ND{F_{\Delta_1}}$ is non-empty, then for all $\gamma\in(0,1)$, the inequality is strict for all $\Diid\in\DSetBCiid$ and $\E{s_1\sim \Delta_1}{\pwr[s_1,\gamma][\Dbd]}\not\geqMost[][\DSetBd] \E{s_2\sim \Delta_2}{\pwr[s_2,\gamma][\Dbd]}$.

These results also hold when replacing $F_{\Delta_i}$ with $F_{\Delta_i}^*\defeq \set{\E{s_i\sim \Delta_i}{\fpi{s_i}}\mid \pi\in\Pi}$ for $i=1,2$.
\end{restatable}
\begin{proof} 
\begin{align}
    \phi\cdot \ND{\NormF{F_{\Delta_1},\gamma}}&\subseteq \phi\cdot \NormF{\ND{F_{\Delta_1}},\gamma}\label{eq:subset-nd-norm}\\
    &\defeq \set{\permute\lim_{\gamma^*\to\gamma}(1-\gamma^*)\f(\gamma^*) \mid \f \in \ND{F_{\Delta_1}}}\\
    &=\set{\lim_{\gamma^*\to\gamma}(1-\gamma^*)\permute\f(\gamma^*) \mid \f \in \ND{F_{\Delta_1}}}\label{eq:cont-permute-app}\\
    &=\set{\lim_{\gamma^*\to\gamma}(1-\gamma^*)\f(\gamma^*) \mid \f \in F_{\text{sub}}'}\\
    &\subseteq \set{\lim_{\gamma^*\to\gamma}(1-\gamma^*)\f(\gamma^*) \mid \f\in F_{\Delta_2}}\label{eq:subset-f}\\
    &\eqdef \NormF{F_{\Delta_2},\gamma}.\label{eq:subset-f-final}
\end{align}
\Cref{eq:subset-nd-norm} follows by \cref{lem:nd-norm-visit}. \Cref{eq:cont-permute-app} follows because $\permute$ is a continuous linear operator. \Cref{eq:subset-f} follows by assumption.
\begin{align}
    \E{s_1\sim \Delta_1}{\pwr[s_1,\gamma]}&\defeq\E{\substack{s_1\sim \Delta_1,\\\rf\sim\Dbd}}{\max_{\pi\in\Pi} \lim_{\gamma^*\to\gamma}\frac{1-\gamma^*}{\gamma^*}\prn{\fpi{s_1}(\gamma^*)-\unitvec[s_1]}^\top \rf}\label{eq:expect-pwr-1}\\
    &=\E{\rf\sim\Dbd}{\max_{\pi\in\Pi} \lim_{\gamma^*\to\gamma}\frac{1-\gamma^*}{\gamma^*}\E{s_1\sim \Delta_1}{\fpi{s_1}(\gamma^*)-\unitvec[s_1]}^\top \rf}\label{eq:bring-inside-expect}\\
    &=\E{\rf\sim\Dbd}{\max_{\dbf\in \NormF{F_{\Delta_1},\gamma}} \dbf^\top \rf}\\
    &=\E{\rf\sim\Dbd}{\max_{\dbf\in \ND{\NormF{F_{\Delta_1},\gamma}}} \dbf^\top \rf}\label{eq:restrict-nd-fdelta}\\
    &\leqMost[][\DSetBd] \E{\rf\sim\Dbd}{\max_{\dbf\in \NormF{F_{\Delta_2},\gamma}} \dbf^\top \rf}\label{eq:leq-most-apply-transient}\\
    &=\E{\rf\sim\Dbd}{\max_{\pi\in\Pi} \lim_{\gamma^*\to\gamma}\frac{1-\gamma^*}{\gamma^*}\E{s_2\sim \Delta_2}{\fpi{s_2}(\gamma^*)-\unitvec[s_2]}^\top \rf}\\
    &=\E{\substack{s_2\sim \Delta_2,\\\rf\sim\Dbd}}{\max_{\pi\in\Pi} \lim_{\gamma^*\to\gamma}\frac{1-\gamma^*}{\gamma^*}\prn{\fpi{s_2}(\gamma^*)-\unitvec[s_2]}^\top \rf}\label{eq:bring-inside-expect-2}\\
    &\eqdef \E{s_2\sim \Delta_2}{\pwr[s_2,\gamma]}.\label{eq:expect-pwr-2}
\end{align}

\Cref{eq:expect-pwr-1} and \cref{eq:expect-pwr-2} follow by \cref{lem:pwr-limit}. \Cref{eq:bring-inside-expect} and \cref{eq:bring-inside-expect-2} follow because each $R$ has a stationary deterministic optimal policy $\pi\in \optPi\subseteq \Pi$ which simultaneously achieves  optimal value at all states. \Cref{eq:restrict-nd-fdelta} follows by \cref{cor:nd-func-indif}.

Apply \cref{lem:expect-superior} with $A\defeq \NormF{F_{\Delta_1},\gamma}, B \defeq \NormF{F_{\Delta_2},\gamma}$, $g$ the identity function, and involution $\phi$ (satisfying $\phi\cdot \ND{A}\subseteq B$ by \cref{eq:subset-f-final}) in order to conclude that \cref{eq:leq-most-apply-transient} holds. 

Suppose that $\ND{F_{\Delta_2}}\setminus \phi\cdot \ND{F_{\Delta_1}}$ is non-empty; let $F_\text{sub}'\defeq \phi\cdot \ND{F_{\Delta_1}}$. \Cref{lem:nd-gamma-F-subset} shows that for all $\gamma\in(0,1)$, $\ND{F_{\Delta_2}(\gamma)}\setminus F_\text{sub}'(\gamma)$ is non-empty. \Cref{lem:pos-aff-nd-invar} \cref{item:invar-vectors} then implies that $\ND{B}\setminus \phi\cdot A=\frac{1-\gamma}{\gamma}\prn{\ND{F_{\Delta_2}(\gamma)}-\unitvec}\setminus \prn{\frac{1-\gamma}{\gamma}F'_\text{sub}(\gamma)}$ is non-empty. Then \cref{lem:expect-superior} implies that for all $\gamma\in(0,1)$, \cref{eq:leq-most-apply-transient} is strict for all $\Diid\in\DSetBCiid$ and $\E{s_1\sim \Delta_1}{\pwr[s_1,\gamma][\Dbd]}\not\geqMost[][\DSetBd] \E{s_2\sim \Delta_2}{\pwr[s_2,\gamma][\Dbd]}$. 

We show that this result's preconditions holding for  $F_{\Delta_i}^*$ implies the $F_{\Delta_i}$ preconditions. Suppose $F_{\Delta_i}^*\defeq \set{\E{s_i\sim \Delta_i}{\fpi{s_i}}\mid \pi\in\Pi}$ for $i=1,2$ are such that $F_\text{sub}^*\defeq \phi\cdot \ND{F_{\Delta_1}^*}\subseteq F_{\Delta_2}^*$. In the following, the $\Delta_i$ are represented as vectors in $\rewardVS$, and $\gamma$ is a variable.
\begin{align}
    \phi\cdot \set{\gamma\f\mid \f \in \ND{F_{\Delta_1}}}&=\phi\cdot \prn{\ND{F_{\Delta_1}^*-\Delta_1}}\\
    &=\phi\cdot\prn{\ND{F_{\Delta_1}^*}-\Delta_1}\label{eq:invariant-vec-fn-set}\\
    &=\set{\permute \f - \permute\Delta_1\mid \f\in \ND{F_{\Delta_1}^*}}\\
    &\subseteq \set{\f-\Delta_2 \mid \f\in F_{\Delta_2}^*}\label{eq:fsub-no-state}\\
    &=\set{\gamma\f\mid \f \in F_{\Delta_2}}.\label{eq:gamma-sim}
\end{align}

\Cref{eq:invariant-vec-fn-set} follows from \cref{lem:pos-aff-nd-invar} \cref{item:invar-vector-fns}. Since we assumed that $\phi\cdot \ND{F_{\Delta_1}^*}\subseteq F_{\Delta_2}^*$, $\phi\cdot \set{\Delta_1}=\phi\cdot\prn{\ND{F_{\Delta_1}^*}(0)}\subseteq F_{\Delta_2}^*(0) =\set{\Delta_2}$. This implies that $\permute\Delta_1=\Delta_2$ and so \cref{eq:fsub-no-state} follows. 

\Cref{eq:gamma-sim} shows that $\phi\cdot \set{\gamma\f\mid \f \in \ND{F_{\Delta_1}}}\subseteq \set{\gamma\f\mid \f \in F_{\Delta_2}}$. But we then have $\phi\cdot \set{\gamma\f\mid \f \in \ND{F_{\Delta_1}}}\defeq \set{\gamma\permute \f \mid \f \in \ND{F_{\Delta_1}}}=\set{\gamma\f \mid \f \in \phi\cdot \ND{F_{\Delta_1}}}\subseteq \set{\gamma\f\mid \f \in F_{\Delta_2}}$. Thus,  $\phi\cdot \ND{F_{\Delta_1}}\subseteq F_{\Delta_2}$. 

Suppose $\ND{F_{\Delta_2}^*}\setminus \phi\cdot \ND{F_{\Delta_1}^*}$ is non-empty, which implies that
\begin{align}
    \phi\cdot \set{\gamma\f\mid \f \in \ND{F_{\Delta_1}}}&=\set{\permute \f - \permute\Delta_1\mid \f\in \ND{F_{\Delta_1}^*}}\\
    &=\set{\f - \permute\Delta_1\mid \f\in \phi\cdot \ND{F_{\Delta_1}^*}}\\
    &\subsetneq \set{\f-\Delta_2 \mid \f\in \ND{F_{\Delta_2}^*}}\\
    &=\set{\gamma\f\mid \f \in \ND{F_{\Delta_2}}}.
\end{align}

Then $\ND{F_{\Delta_2}}\setminus \phi\cdot \ND{F_{\Delta_1}}$ must be non-empty. Therefore, if the preconditions of this result are met for $F_{\Delta_i}^*$, they are met for $F_{\Delta_i}$.
\end{proof}

\morePowerMoreOptions*
\begin{proof}
Let $F_\text{sub}\defeq \phi\cdot \Fnd(s')\subseteq \F(s)$. Let $\Delta_1\defeq \unitvec[s'],\Delta_2\defeq \unitvec$, and define $F_{\Delta_i}^*\defeq \set{\E{s_i\sim \Delta_i}{\fpi{s_i}}\mid \pi\in\Pi}$ for $i=1,2$. Then $\Fnd(s')=\ND{F_{\Delta_1}^*}$ is similar to $F_\text{sub}=F^*_\text{sub}\subseteq F_{\Delta_2}^*=\F(s)$ via involution $\phi$. Apply \cref{lem:more-power-prefix} to conclude that  $\forall \gamma\in[0,1]:\pwr[s',\gamma][\Dbd]\leqMost[][\DSetBd] \pwr[s,\gamma][\Dbd]$. 

Furthermore, $\Fnd(s)=\ND{F_{\Delta_2}^*}$, and $F_\text{sub}=F_\text{sub}^*$, and so if $\Fnd(s)\setminus \phi\cdot \Fnd(s')\defeq \Fnd(s)\setminus F_\text{sub} =\ND{F_{\Delta_2}^*}\setminus F_\text{sub}^*$ is non-empty, then \cref{lem:more-power-prefix} shows that for all $\gamma\in(0,1)$, the inequality is strict for all $\Diid\in\DSetBCiid$ and $\pwr[s',\gamma][\Dbd]\not\geqMost[][\DSetBd] \pwr[s,\gamma][\Dbd]$. 
\end{proof}

\begin{restatable}[Non-dominated visit distribution functions never agree with other visit distribution functions at that state]{lem}{noAgreeND}\label{lem:no-agree} 
Let $\f \in \Fnd(s), \f'\in\F(s)\setminus\{\f\}$. $\forall \gamma \in (0,1):\f(\gamma) \neq \f'(\gamma)$.
\end{restatable}
\begin{proof}
Let $\gamma\in(0,1)$. Since $\f\in\Fnd(s)$, there exists a $\gamma^*\in(0,1)$ at which $\f$ is strictly optimal for some reward function. Then by \cref{transferDiscount}, we can produce another reward function for which $\f$ is strictly optimal at discount rate $\gamma$; in particular, \cref{transferDiscount} guarantees that the policies which induce $\f'$ are not optimal at $\gamma$. So $\f(\gamma)\neq \f'(\gamma)$. 
\end{proof}

\begin{restatable}[Cardinality of non-dominated visit distributions]{cor}{cardNDInter}\label{cor:card-nd-visit}
Let $F\subseteq \F(s)$. $\forall \gamma\in(0,1): \abs{F\cap \Fnd(s)}=\abs{F(\gamma)\cap \Fnd(s,\gamma)}$.
\end{restatable}
\begin{proof}
Let $\gamma\in(0,1)$. By applying \cref{lem:nd-gamma-F-subset} with $\Delta_d\defeq \unitvec$, $\f\in\Fnd(s)=\ND{\F(s)}$ iff $\f(\gamma)\in\ND{\F(s,\gamma)}$. By \cref{lem:nd-relation}, $\ND{\F(s,\gamma)}=\Fnd(s,\gamma)$. So all $\f\in F\cap \Fnd(s)$ induce $\f(\gamma)\in F(\gamma)\cap \Fnd(s,\gamma)$, and $\abs{F\cap \Fnd(s)}\geq\abs{F(\gamma)\cap \Fnd(s,\gamma)}$.

\Cref{lem:no-agree} implies that for all $\f,\f'\in\Fnd(s)$, $\f=\f'$ iff $\f(\gamma)=\f'(\gamma)$. Therefore, $\abs{F\cap \Fnd(s)}\leq\abs{F(\gamma)\cap \Fnd(s,\gamma)}$. So $\abs{F\cap \Fnd(s)}=\abs{F(\gamma)\cap \Fnd(s,\gamma)}$.
\end{proof}

\begin{restatable}[Optimality probability and state bottlenecks]{lem}{optProbBottle}\label{lem:opt-prob-bottleneck}
Suppose that $s$ can reach $\reach{s',a'}\cup\reach{s',a}$, but only by taking actions equivalent to $a'$ or $a$ at state $s'$. $F_{\text{nd},a'}\defeq \FndRestrictAction{s'}{a'}, F_a\defeq \FRestrictAction{s'}{a}$. Suppose $F_a$ contains a copy of $F_{\text{nd},a'}$ via $\phi$ which fixes all states not belonging to $\reach{s',a'}\cup\reach{s',a}$. 
Then $\forall \gamma\in[0,1]:\optprob[\Dany]{F_{\text{nd},a'},\gamma}\leqMost[][\DSetAny] \optprob[\Dany]{F_a,\gamma}$.

If $\Fnd(s)\cap \prn{F_a\setminus \phi\cdot F_{\text{nd},a'}}$ is non-empty, then for all $\gamma\in(0,1)$, the inequality is strict for all $\Diid\in\DSetBCiid$, and $\optprob[\Dany]{F_{\text{nd},a'},\gamma}\not\geqMost[][\DSetAny] \optprob[\Dany]{F_a,\gamma}$.
\end{restatable}
\begin{proof} Let $F_\text{sub}\defeq \phi\cdot F_{\text{nd},a'}$. Let $F^* \defeq \bigcup_{\substack{a''\in\A:\\ \prn{a'' \not\equiv_{s'} a} \land \prn{a'' \not \equiv_{s'}a'}}} \FRestrictAction{s'}{a''}\cup F_{\text{nd},a'} \cup F_\text{sub}$.
\begin{align}
    \phi\cdot F^*\defeq\,& \phi\cdot \prn{\bigcup_{\substack{a''\in\A:\\ \prn{a'' \not\equiv_{s'} a} \land \prn{a'' \not \equiv_{s'}a'}}} \FRestrictAction{s'}{a''}\cup F_{\text{nd},a'} \cup F_\text{sub}}\\
    =\,& \bigcup_{\substack{a''\in\A:\\ \prn{a'' \not\equiv_{s'} a} \land \prn{a'' \not \equiv_{s'}a'}}}\phi\cdot \FRestrictAction{s'}{a''}\cup \prn{\phi\cdot F_{\text{nd},a'}} \cup \prn{\phi\cdot F_\text{sub}}\\
    =\,& \bigcup_{\substack{a''\in\A:\\ \prn{a'' \not\equiv_{s'} a} \land \prn{a'' \not \equiv_{s'}a'}}}\phi\cdot \FRestrictAction{s'}{a''}\cup F_\text{sub} \cup F_{\text{nd},a'}\label{eq:invert-Fsub}\\
    =\,& \bigcup_{\substack{a''\in\A:\\ \prn{a'' \not\equiv_{s'} a} \land \prn{a'' \not \equiv_{s'}a'}}} \FRestrictAction{s'}{a''}\cup F_\text{sub}\cup F_{\text{nd},a'}\label{eq:phi-prime-bigcup}\\
    \eqdef\,& F^*.
\end{align}
\Cref{eq:invert-Fsub} follows because the involution $\phi$ ensures that $\phi\cdot F_\text{sub}=F_{\text{nd},a'}$. By assumption, $\phi$ fixes all $s'\not\in \reach{s',a'}\cup\reach{s',a}$. Suppose $\f\in\F (s)\setminus\prn{F_{\text{nd},a'}\cup F_{a}}$. By the bottleneck assumption, $\f$ does not visit states in $\reach{s',a'}\cup\reach{s',a}$. Therefore, $\permute[\phi]\f=\f$, and so \cref{eq:phi-prime-bigcup} follows.

Let $F_Z \defeq \prn{\F(s)\setminus (\FRestrictAction{s}{a'}\cup F_a)} \cup F_{\text{nd},a'} \cup F_a$. By definition, $F_Z\subseteq \F(s)$. Furthermore, $\Fnd(s)=\bigcup_{\substack{a''\in\A}} \FndRestrictAction{s'}{a''}\subseteq \prn{\F(s)\setminus (\FRestrictAction{s}{a'}\cup F_a)} \cup \FndRestrictAction{s}{a'} \cup F_a\eqdef F_Z$, and so $\Fnd(s)\subseteq F_Z$. Note that $F^*=F_Z\setminus (F_a\setminus F_\text{sub})$.

\paragraph*{Case: $\gamma\in(0,1)$.}
\begin{align}
    \optprob[\Dany]{F_{\text{nd},a'},\gamma} &= \phelper{F_{\text{nd},a'}(\gamma)\geq \F(s,\gamma)}[\Dany]\label{eq:opt-prb-id-1}\\
    &\leqMost[][\DSetAny] \phelper{F_{a}(\gamma)\geq \F(s,\gamma)}[\Dany]\label{eq:leq-most-apply}\\
    &= \optprob[\Dany]{F_{\text{nd},a'},\gamma}.\label{eq:opt-prb-id-2}
\end{align} 

\Cref{eq:opt-prb-id-1} and \cref{eq:opt-prb-id-2} follow from \cref{lem:opt-prob-id}. \Cref{eq:leq-most-apply} follows by applying \cref{lem:opt-prob-superior} with $A\defeq F_{\text{nd},a'}(\gamma),B'\defeq F_\text{sub}(\gamma),B\defeq F_{a}(\gamma),C\defeq\F(s,\gamma),Z\defeq F_Z(\gamma)$ which satisfies $\ND{C}=\Fnd(s,\gamma)\subseteq F_Z(\gamma)\subseteq \F(s,\gamma)=C$, and involution $\phi$ which satisfies $\phi\cdot F^*(\gamma)=\phi\cdot \prn{Z\setminus\prn{B\setminus B'}}=Z\setminus\prn{B\setminus B'}=F^*(\gamma)$. 

Suppose $\Fnd(s)\cap \prn{F_a\setminus F_\text{sub}}$ is non-empty. $0<\abs{\Fnd(s)\cap \prn{F_a\setminus F_\text{sub}}}=\abs{\Fnd(s,\gamma)\cap\prn{F_a(\gamma)\setminus F_\text{sub}(\gamma)}}\eqdef\abs{\ND{C}\cap\prn{B\setminus B'}}$ (with the first equality holding by \cref{cor:card-nd-visit}), and so $\ND{C}\cap\prn{B\setminus B'}$ is non-empty. We also have $B\defeq F_{a}(\gamma)\subseteq \F(s,\gamma)\eqdef C$. Then reapplying \cref{lem:opt-prob-superior}, \cref{eq:leq-most-apply} is strict for all $\Diid\in\DSetBCiid$, and $\optprob[\Dany]{F_{\text{nd},a'},\gamma}\not\geqMost[][\DSetAny] \optprob[\Dany]{F_a,\gamma}$.

\paragraph*{Case: $\gamma=1$, $\gamma=0$.} 
\begin{align}
    \optprob[\Dany]{F_{\text{nd},a'},1} &=\lim_{\gamma^*\to 1} \optprob[\Dany]{F_{\text{nd},a'},\gamma^*}\label{eq:lim-1-prob}\\ 
    &=\lim_{\gamma^*\to 1} \phelper{F_{\text{nd},a'}(\gamma^*)\geq \F(s,\gamma^*)}[\Dany]\label{eq:lim-1-opt-prb-id-1}\\
    &\leqMost[][\DSetAny] \lim_{\gamma^*\to 1} \phelper{F_{a}(\gamma^*)\geq \F(s,\gamma^*)}[\Dany]\label{eq:apply-lim-prob}\\
    &=\lim_{\gamma^*\to 1} \optprob[\Dany]{F_{a},\gamma^*}\label{eq:lim-1-opt-prb-id-2}\\ 
    &=\optprob[\Dany]{F_{a},1}.\label{eq:lim-1-prob-b}
\end{align}
\Cref{eq:lim-1-prob} and \cref{eq:lim-1-prob-b} hold by \cref{prop:opt-prob-converge}. \Cref{eq:lim-1-opt-prb-id-1} and \cref{eq:lim-1-opt-prb-id-2} follow by \cref{lem:opt-prob-id}. Applying \cref{lem:lim-prob-most} with $\gamma\defeq 1,I\defeq (0,1), F_A\defeq F_{\text{nd},a'}, F_B\defeq F_a, F_C\defeq \F(s)$, $F_Z$ as defined above, and involution $\phi$ (for which $\phi\cdot \prn{F_Z\setminus \prn{F_B\setminus \phi\cdot F_A}}=F_Z\setminus \prn{F_B\setminus \phi\cdot F_A}$),  we conclude that \cref{eq:apply-lim-prob} follows.

The $\gamma=0$ case proceeds similarly to $\gamma=1$.
\end{proof}

\begin{restatable}[Action optimality probability is a special case of visit distribution optimality probability]{lem}{agreeOptProb}\label{lem:agree-opt-prob} $\optprob[\Dany]{s,a,\gamma}=\optprob[\Dany]{\FRestrictAction{s}{a},\gamma}$.
\end{restatable}
\begin{proof}
Let $F_a\defeq \FRestrictAction{s}{a}$. For $\gamma \in (0,1)$,
\begin{align}
    \prob[\Dany]{s,a,\gamma}&\defeq\prob[R \sim \Dany]{\exists \pi^* \in \optPi: \pi^*(s)=a}\\
    &= \prob[\rf \sim \Dany]{\exists \fpi[\pi^*]{s} \in F_a: \fpi[\pi^*]{s}(\gamma)^\top \rf=\max_{\f \in \F(s)}\f(\gamma)^\top \rf}\label{eq:f-action-equiv}\\
    &=\optprob[\Dany]{F_a,\gamma}.\label{eq:f-action-equiv-final}
\end{align}
By \cref{lem:opt-pol-visit-iff}, if $\exists\pi^*\in\optPi:\pi^*(s)=a$, then it induces some optimal $\fpi[\pi^*]{s}\in F_a$. Conversely, if $\fpi[\pi^*]{s}\in F_a$ is optimal at $\gamma\in(0,1)$, then $\pi^*$ chooses optimal actions on the support of $\fpi[\pi^*]{s}(\gamma)$. Let $\pi'$ agree with $\pi^*$ on that support and let $\pi'$ take optimal actions at all other states. Then $\pi'\in\optPi$ and $\pi'(s)=a$. So \cref{eq:f-action-equiv} follows.

Suppose $\gamma=0$ or $\gamma=1$. Consider any sequence $\prn{\gamma_n}_{n=1}^\infty$ converging to $\gamma$, and let $\Dany$ induce probability measure $F$.
\begin{align}
    \optprob[\Dany]{F_a,\gamma}&\defeq \lim_{\gamma^*\to \gamma} \optprob[\Dany]{F_a,\gamma^*}\\
    &=\lim_{\gamma^*\to \gamma}\prob[R \sim \Dany]{\exists \pi^* \in \optPi[R,\gamma^*]: \pi^*(s)=a}\label{eq:inner-eq-act}\\
    &=\lim_{n\to\infty}\prob[R \sim \Dany]{\exists \pi^* \in \optPi[R,\gamma_n]: \pi^*(s)=a}\\
    &=\lim_{n\to\infty}\int_{\rewardSpace} \indic{\exists \pi^* \in \optPi[R,\gamma_n]: \pi^*(s)=a} \dF[R]\\
    &=\int_{\rewardSpace}\lim_{n\to\infty} \indic{\exists \pi^* \in \optPi[R,\gamma_n]: \pi^*(s)=a} \dF[R]\label{eq:dominated-convergence}\\
    &=\int_{\rewardSpace} \indic{\exists \pi^* \in \optPi[R,\gamma]: \pi^*(s)=a} \dF[R]\\
    &\eqdef \prob[\Dany]{s,a,\gamma}.
\end{align}

\Cref{eq:inner-eq-act} follows by \cref{eq:f-action-equiv-final}. for $\gamma^*\in[0,1]$, let $f_{\gamma^*}(R)\defeq \indic{\exists \pi^* \in \optPi[R,\gamma^*]: \pi^*(s)=a}$. For each $R\in\rewardSpace$, \cref{lem:opt-pol-shift-bound} exists $\gamma_x\approx \gamma$ such that for all intermediate $\gamma_x'$ between $\gamma_x$ and $\gamma$, $ \optPi[R,\gamma_x']=\optPi[R,\gamma]$. Since $\gamma_n \to \gamma$, this means that $\prn{f_{\gamma_n}}_{n=1}^\infty$ converges pointwise to $f_\gamma$. Furthermore, $\forall n \in \mathbb{N}, R\in\rewardSpace: \abs{f_{\gamma_n}(R)}\leq 1$ by definition. Therefore, \cref{eq:dominated-convergence} follows by Lebesgue's dominated convergence theorem.
\end{proof}

\graphOptions*
\begin{proof}
Note that by \cref{def:visit}, $F_{a'}(0)=\set{\unitvec}=F_a(0)$. Since $\phi\cdot F_{a'}\subseteq F_a$, in particular we have $\phi\cdot F_{a'}(0)=\set{\permute\unitvec}\subseteq \set{\unitvec}=F_a(0)$, and so $\phi(s)=s$.

\textbf{\Cref{item:power-options}.}
For state probability distribution $\Delta_s\in \Delta(\St)$, let $F^*_{\Delta_s}\defeq \set{\E{s'\sim \Delta_s}{\fpi{s'}}\mid \pi\in\Pi}$. Unless otherwise stated, we treat $\gamma$ as a variable in this item; we apply element-wise vector addition, constant multiplication, and variable multiplication via the conventions outlined in \cref{def:aff-transf-short}.
\begin{align}
    F_{a'} &= \set{\unitvec + \gamma \E{s_{a'}\sim T(s,a')}{\fpi{s_{a'}}}\mid \pi\in\Pi: \pi(s)=a'}\label{eq:visit-defn}\\
    &= \set{\unitvec + \gamma \E{s_{a'}\sim T(s,a')}{\fpi{s_{a'}}}\mid \pi\in\Pi}\label{eq:remove-action-constraint}\\
    &= \unitvec+\gamma F^*_{T(s,a')}.\label{eq:Fstar-contain}
\end{align}

\Cref{eq:visit-defn} follows by \cref{def:visit}, since each $\f\in\F(s)$ has an initial term of $\unitvec$. \Cref{eq:remove-action-constraint} follows because $s\not\in \reach{s,a'}$, and so for all $s_{a'}\in \supp[T(s,a')]$, $\fpi{s_{a'}}$ is unaffected by the choice of action $\pi(s)$. Note that similar reasoning implies that $F_a\subseteq \unitvec+\gamma F^*_{T(s,a)}$ (because \cref{eq:remove-action-constraint} is a containment relation in general).

Since $F_{a'}=\unitvec+\gamma F^*_{T(s,a')}$, if $F_a$ contains a copy of $F_{a'}$ via $\phi$, then $F^*_{T(s,a)}$ contains a copy of $F^*_{T(s,a')}$ via $\phi$. Then $\phi\cdot \ND{F^*_{T(s,a')}}\subseteq \phi\cdot F^*_{T(s,a')}\subseteq F^*_{T(s,a)}$, and so $F^*_{T(s,a)}$ contains a copy of $\ND{F^*_{T(s,a')}}$. Then apply \cref{lem:more-power-prefix} with $\Delta_1\defeq T(s,a')$ and $\Delta_2\defeq T(s,a)$ to conclude that $\forall\gamma\in[0,1]:\E{s_{a'}\sim T(s,a')}{\pwr[s_{a'},\gamma][\Dbd]}\leqMost[][\DSetBd] \E{s_{a}\sim T(s,a)}{\pwr[s_a,\gamma][\Dbd]}$. 

Suppose $\Fnd(s)\cap \prn{F_a\setminus \phi\cdot F_{a'}}$ is non-empty. To apply the second condition of \cref{lem:more-power-prefix}, we want to demonstrate that $\ND{F^*_{T(s,a)}}\setminus \phi\cdot \ND{F^*_{T(s,a')}}$ is also non-empty.

First consider $\f\in\Fnd(s)\cap F_a$. Because $F_a\subseteq \unitvec+\gamma F^*_{T(s,a)}$, we have that $\gamma\inv(\f - \unitvec)\in F^*_{T(s,a)}$. Because $\f \in \Fnd(s)$, by \cref{def:nd}, $\exists \rf \in \rewardVS, \gamma_x\in (0,1)$ such that
\begin{align}
    \f(\gamma_x)^\top\rf  &> \max_{\f'\in\F(s)\setminus \set{\f}}\f'(\gamma_x)^\top \rf.\label{eq:f-nd-def}
\end{align}
Then since $\gamma_x\in (0,1)$,
\begin{align}
    \gamma_x\inv(\f(\gamma_x)-\unitvec)^\top\rf  &> \max_{\f'\in\F(s)\setminus \set{\f}}\gamma_x\inv(\f'(\gamma_x)-\unitvec)^\top \rf\\
    &=\max_{\f'\in\gamma_x\inv\prn{(\F(s)\setminus \set{\f})-\unitvec}}\f'(\gamma_x)^\top \rf\\
    &\geq \max_{\f'\in \gamma_x\inv\prn{(F_a\setminus \set{\f})-\unitvec}}\f'(\gamma_x)^\top \rf\label{eq:Fact-contain}\\
    &= \max_{\f'\in F^*_{T(s,a)}\setminus \set{\gamma_x\inv(\f-\unitvec)}}\f'(\gamma_x)^\top \rf.\label{eq:FTsa}
\end{align}

\Cref{eq:Fact-contain} holds because $F_a\subseteq \F(s)$. By assumption, action $a$ is optimal for $\rf$ at state $s$ and at discount rate $\gamma_x$. \Cref{eq:remove-action-constraint} shows that $F^*_{T(s,a)}$ potentially allows the agent a non-stationary policy choice at $s$, but  non-stationary policies cannot increase optimal value \citep{puterman_markov_2014}. Therefore, \cref{eq:FTsa} holds.

We assumed that $\gamma\inv (\f-\unitvec)\in \gamma\inv(\Fnd(s)-\unitvec)$. Furthermore, since we just showed that $\gamma\inv (\f-\unitvec)\in F^*_{T(s,a)}$ is strictly optimal over the other elements of $F^*_{T(s,a)}$ for reward function $\rf$ at discount rate $\gamma_x\in(0,1)$, we conclude that it is an element of $\ND{F^*_{T(s,a)}}$ by \cref{def:nd-vec-func}. Then we conclude that $\gamma\inv(\Fnd(s)-\unitvec)\cap F^*_{T(s,a)}\subseteq \ND{F^*_{T(s,a)}}$. 

We now show that $\ND{F^*_{T(s,a)}}\setminus \phi\cdot \ND{F^*_{T(s,a')}}$ is non-empty. 
\begin{align}
    0&<\abs{\Fnd(s)\cap \prn{F_a\setminus \phi\cdot F_{a'}}}\label{eq:nonempty-fnd-cap}\\
    &=\abs{\gamma\inv\prn{\Fnd(s)\cap \prn{F_a\setminus \phi\cdot F_{a'}}-\unitvec}}\label{eq:card-preserve}\\
    &\leq\abs{\gamma\inv\prn{\Fnd(s)-\unitvec}\cap \prn{F^*_{T(s,a)}\setminus\phi\cdot F^*_{T(s,a')}}}\label{eq:Fstar-convert}\\
    &=\abs{\prn{\gamma\inv\prn{\Fnd(s)-\unitvec}\cap F^*_{T(s,a)}}\setminus\phi\cdot F^*_{T(s,a')}}\\
    &\leq\abs{\ND{F^*_{T(s,a)}}\setminus\phi\cdot F^*_{T(s,a')}}\label{eq:contain-nd-fstar}\\
    &\leq\abs{\ND{F^*_{T(s,a)}}\setminus \phi\cdot \ND{F^*_{T(s,a')}}}.\label{eq:nd-containment-trivial}
\end{align}
\Cref{eq:nonempty-fnd-cap} follows by the assumption that $\Fnd(s)\cap \prn{F_a\setminus \phi\cdot F_{a'}}$ is non-empty. Let $\f,\f'\in \Fnd(s)\cap \prn{F_a\setminus \phi\cdot F_{a'}}$ be distinct. Then we must have that for some $\gamma_x\in (0,1)$, $\f(\gamma_x)\neq \f'(\gamma_x)$. This holds iff $\gamma_x\inv(\f(\gamma_x)-\unitvec)\neq \gamma_x\inv(\f'(\gamma_x)-\unitvec)$, and so \cref{eq:card-preserve} holds.

\Cref{eq:Fstar-convert} holds because $F_a\subseteq \unitvec+\gamma F^*_{T(s,a)}$ and $F_a'= \unitvec+\gamma F^*_{T(s,a')}$ by \cref{eq:Fstar-contain}. \Cref{eq:contain-nd-fstar} holds because we showed above that $\gamma\inv(\Fnd(s)-\unitvec)\cap F^*_{T(s,a)}\subseteq \ND{F^*_{T(s,a)}}$. \Cref{eq:nd-containment-trivial} holds because $\ND{F^*_{T(s,a')}}\subseteq F^*_{T(s,a')}$ by \cref{def:nd-vec-func}. 

Therefore, $\ND{F^*_{T(s,a)}}\setminus \phi\cdot \ND{F^*_{T(s,a')}}$ is non-empty, and so apply the second condition of \cref{lem:more-power-prefix} to conclude that for all $\Diid\in\DSetBCiid$, $\forall\gamma\in(0,1):\E{s_{a'}\sim T(s,a')}{\pwr[s_{a'},\gamma][\Diid]}< \E{s_{a}\sim T(s,a)}{\pwr[s_a,\gamma][\Diid]}$, and that $\forall\gamma\in(0,1):\E{s_{a'}\sim T(s,a')}{\pwr[s_{a'},\gamma][\Dbd]}\not \geqMost[][\DSetBd] \E{s_{a}\sim T(s,a)}{\pwr[s_a,\gamma][\Dbd]}$.

\textbf{\Cref{item:opt-prob-options}.} 
Let $\phi'(s_x)\defeq \phi(s_x)$ when $s_x\in\reach{s,a'}\cup\reach{s,a}$, and equal $s_x$ otherwise. Since $\phi$ is an involution, so is $\phi'$. 
\begin{align}
    \phi'\cdot F_{a'}&\defeq \set{\permute[\phi']\prn{\unitvec + \gamma\E{s_{a'}\sim T(s,a')}{ \fpi{s_{a'}}}}\mid \pi\in \Pi, \pi(s)=a'}\\
    &= \set{\unitvec+\gamma\E{s_{a'}\sim T(s,a')}{\permute[\phi'] \fpi{s_{a'}}}\mid \pi\in \Pi, \pi(s)=a'}\label{eq:fix-unitvec}\\
    &= \set{\permute\unitvec+ \gamma\E{s_{a'}\sim T(s,a')}{\permute[\phi] \fpi{s_{a'}}}\mid \pi\in \Pi, \pi(s)=a'}\label{eq:phi-prime-eq}\\
    &\eqdef \phi\cdot F_{a'}\\
    &\subseteq F_a.\label{eq:sim-subtracted}
\end{align}

\Cref{eq:fix-unitvec} follows because if $s\in\reach{s,a'}\cup\reach{s,a}$, then we already showed that $\phi$ fixes $s$. Otherwise, $\phi'(s)=s$ by definition. \Cref{eq:phi-prime-eq} follows by the definition of $\phi'$ on $\reach{s,a'}\cup\reach{s,a}$ and because $\unitvec=\permute\unitvec$. Next, we assumed that $\phi\cdot F_{a'}\subseteq F_a$, and so \cref{eq:sim-subtracted} holds.

Therefore, $F_a$ contains a copy of $F_{a'}$ via $\phi'$ fixing all $s_x\not\in\reach{s,a'}\cup\reach{s,a}$. Therefore, $F_a$ contains a copy of $F_{\text{nd},a'}\defeq \Fnd(s)\cap F_{a'}$ via the same $\phi'$. Then apply \cref{lem:opt-prob-bottleneck} with $s'\defeq s$ to conclude that $\forall \gamma\in[0,1]:\optprob[\Dany]{F_{a'},\gamma}\leqMost[][\DSetAny] \optprob[\Dany]{F_a,\gamma}$. By \cref{lem:agree-opt-prob}, $\optprob[\Dany]{s,a',\gamma}=\optprob[\Dany]{F_{a'},\gamma}$ and $\optprob[\Dany]{s,a,\gamma}=\optprob[\Dany]{F_a,\gamma}$. Therefore, $\forall \gamma\in[0,1]:\optprob[\Dany]{s,a',\gamma}\leqMost[][\DSetAny] \optprob[\Dany]{s,a,\gamma}$.

If $\Fnd(s)\cap \prn{F_a\setminus \phi\cdot F_{a'}}$ is non-empty, then apply the second condition of \cref{lem:opt-prob-bottleneck} to conclude that for all $\gamma\in(0,1)$, the inequality is strict for all $\Diid\in\DSetBCiid$, and $\optprob[\Dany]{s,a',\gamma}\not\geqMost[][\DSetAny] \optprob[\Dany]{s,a,\gamma}$.
\end{proof}

\subsubsection{When \texorpdfstring{$\gamma=1$}{reward is undiscounted}, optimal policies tend to navigate towards ``larger'' sets of cycles} \label{sec:preservation}

\begin{restatable}[$\pwrNoDist$ identity when $\gamma=1$]{lem}{gammaOnePower}\label{lem:gamma-1-power}
\begin{equation}
\pwr[s,1][\Dbd]=\E{\rf \sim \Dbd}{\max_{\dbf \in \RSD} \dbf^\top \rf}=\E{\rf \sim \Dbd}{\max_{\dbf \in \RSDnd} \dbf^\top \rf}.
\end{equation}
\end{restatable}
\begin{proof}
\begin{align}
    \pwr[s,1][\Dbd]&= \E{\rf\sim\Dbd}{\max_{\fpi{s}\in \F(s)} \lim_{\gamma\to 1} \frac{1-\gamma}{\gamma}\prn{\fpi{s}(\gamma)-\unitvec}^\top \rf} \label{eq:pwr-lim-rsd}\\
 &=\E{\rf \sim \Dbd}{\max_{\dbf \in \RSD} \dbf^\top \rf}\label{eq:def-rsd}\\
 &=\E{\rf \sim \Dbd}{\max_{\dbf \in \RSDnd} \dbf^\top \rf}.\label{eq:nd-restrict-rsd}
\end{align}
\Cref{eq:pwr-lim-rsd} follows by \cref{lem:pwr-limit}. \Cref{eq:def-rsd} follows by the definition of $\RSD$ (\cref{def:rsd}). \Cref{eq:nd-restrict-rsd} follows because for all $\rf \in \rewardVS$, \cref{cor:nd-func-indif} shows that $\max_{\dbf \in \RSD} \dbf^\top \rf=\max_{\dbf \in \ND{\RSD}} \dbf^\top \rf \eqdef\max_{\dbf \in \RSDnd} \dbf^\top \rf$.
\end{proof}

\RSDSimPower* 
\begin{proof} Suppose $\RSDnd[s']$ is similar to $D\subseteq\RSD$ via involution $\phi$. 
\begin{align}
    \pwr[s',1]&=\E{\rf \sim \Dbd}{\max_{\dbf \in \RSDnd[s']} \dbf^\top \rf}\label{eq:pwr-id-rsd}\\
    &\leqMost[][\DSetBd]\E{\rf \sim \Dbd}{\max_{\dbf \in \RSDnd} \dbf^\top \rf}\label{eq:rsd-func-incr}
    \\
    &=\pwr[s,1][\Dbd]\label{eq:pwr-id-rsd-2}
\end{align} 
\Cref{eq:pwr-id-rsd} and \cref{eq:pwr-id-rsd-2} follow from \cref{lem:gamma-1-power}. By applying \cref{lem:expect-superior} with $A\defeq \RSD[s'],B'\defeq D, B\defeq\RSD$ and $g$ the identity function, \cref{eq:rsd-func-incr} follows. 

Suppose $\RSDnd\setminus D$ is non-empty. By the same result, \cref{eq:rsd-func-incr} is a strict inequality for all $\Diid\in\DSetBCiid$, and we conclude that $\pwr[s',1][\Dbd]\not \geqMost[][\DSetBd] \pwr[s,1][\Dbd]$. 
\end{proof}

\rsdIC*
\begin{proof}
Let $D_\text{sub}\defeq\phi\cdot D'$, where $D_\text{sub}\subseteq D$ by assumption. Let $X\defeq \set{s_i\in \St \mid \max_{\dbf \in D'\cup D} \dbf^\top \unitvec[s_i]>0}$. Define
\begin{equation}
    \phi'(s_i)\defeq 
    \begin{cases}
    \phi(s_i) & \text{ if } s_i\in X\label{eq:phi-prime-def}\\
    s_i &\text{ else}.
    \end{cases}
\end{equation}

Since $\phi$ is an involution, $\phi'$ is also an involution. Furthermore, by the definition of $X$, $\phi'\cdot D'=D_\text{sub}$ and $\phi'\cdot D_\text{sub}=D'$ (because we assumed that both equalities hold for $\phi$). 

Let $D^*\defeq D'\cup D_\text{sub} \cup \prn{\RSDnd\setminus (D' \cup D)}$.
\begin{align}
    \phi'\cdot D^*&\defeq \phi'\cdot \prn{D'\cup D_\text{sub} \cup \prn{\RSDnd\setminus (D' \cup D)}} \\
    &=\prn{\phi'\cdot D'}\cup \prn{\phi'\cdot D_\text{sub}} \cup \phi'\cdot\prn{\RSDnd\setminus (D' \cup D)} \\
    &= D_\text{sub} \cup D' \cup \prn{\RSDnd\setminus (D' \cup D)}\label{eq:sim-phi-prime-rsd}\\
    &\eqdef D^*.\label{eq:fixing-rsd-perm}
\end{align}

In \cref{eq:sim-phi-prime-rsd}, we know that $\phi'\cdot D'=D_\text{sub}$ and $\phi'\cdot D_\text{sub}=D'$. We just need to show that $\phi'\cdot \prn{\RSDnd\setminus (D' \cup D)}=\RSDnd\setminus (D' \cup D)$. 

Suppose $\exists s_i\in X, \dbf' \in \RSDnd\setminus (D' \cup D):\dbf'^\top \unitvec[s_i]>0$. By the definition of $X$, $\exists \dbf \in D' \cup D:\dbf^\top \unitvec[s_i]>0$. Then 
\begin{align}
    \dbf^\top \dbf'&=\sum_{j=1}^{\abs{\St}} \dbf^\top (\dbf'\odot\unitvec[s_j])\label{eq:dot-had}\\
    &\geq \dbf^\top (\dbf'\odot\unitvec[s_i])\label{eq:non-neg-rsd-entry}\\
    &=\dbf^\top \prn{(\dbf'^\top \unitvec[s_i])\unitvec[s_i]}\\
    &= (\dbf'^\top \unitvec[s_i]) \cdot (\dbf^\top\unitvec[s_i])\\
    &>0.\label{eq:pos-rsd-agree}
\end{align}

\Cref{eq:dot-had} follows from the definitions of the dot and Hadamard products. \Cref{eq:non-neg-rsd-entry} follows because $\dbf$ and $\dbf'$ have non-negative entries. \Cref{eq:pos-rsd-agree} follows because $\dbf^\top \unitvec[s_i]$ and $\dbf'^\top \unitvec[s_i]$ are both positive. But \cref{eq:pos-rsd-agree} shows that $\dbf^\top\dbf'>0$, contradicting our assumption that $\dbf$ and $\dbf'$ are orthogonal. 

Therefore, such an $s_i$ cannot exist, and $X'\defeq\set{s_i' \in \St \mid \max_{\dbf' \in \RSDnd\setminus (D' \cup D)} \dbf'^\top \unitvec[s_i]>0}\subseteq (\St \setminus X)$. By \cref{eq:phi-prime-def}, $\forall s_i'\in X':\phi'(s_i')=s_i'$. Thus, $\phi'\cdot \prn{\RSDnd\setminus (D' \cup D)}=\RSDnd\setminus (D' \cup D)$, and \cref{eq:sim-phi-prime-rsd} follows. We conclude that $\phi'\cdot D^*=D^*$.

Consider $Z\defeq \prn{\RSDnd\setminus (D' \cup D)} \cup D \cup D'$. First, $Z\subseteq \RSD$ by definition. Second, $\RSDnd = \RSDnd\setminus (D' \cup D) \cup (\RSDnd\cap D') \cup (\RSDnd \cap D) \subseteq Z$. Note that $D^*=Z\setminus (D \setminus D_\text{sub})$.
\begin{align} 
\avgprob[\Dany]{D'}&= \phelper{D'\geq \RSD}[\Dany]\\ 
&\leqMost[][\DSetAny] \phelper{D\geq \RSD}[\Dany]\label{eq:leq-most-1}\\
&= \avgprob[\Dany]{D}.
\end{align} 

Since $\phi\cdot D'\subseteq D$ and $\ND{D'}\subseteq D'$, $\phi\cdot \ND{D'}\subseteq D$. Then \cref{eq:leq-most-1} holds by applying \cref{lem:opt-prob-superior} with $A\defeq D', B'\defeq D_\text{sub}, B\defeq D, C\defeq \RSD$, and the previously defined $Z$ which we showed satisfies $\ND{C}\subseteq Z\subseteq C$. Furthermore, involution $\phi'$ satisfies $\phi'\cdot  B^* = \phi'\cdot \prn{Z\setminus (B \setminus B')}=Z\setminus (B \setminus B') = B^*$ by \cref{eq:fixing-rsd-perm}.

When $\RSDnd\cap\prn{D\setminus D_\text{sub}}$ is non-empty, since $B'\subseteq C$ by assumption, \cref{lem:opt-prob-superior} also shows that \cref{eq:leq-most-1} is strict for all $\Diid\in\DSetBCiid$, and that $\avgprob[\Dany]{D'}\not\geqMost[][\DSetAny] \avgprob[\Dany]{D}$.
\end{proof} 

\begin{restatable}[{\rsd[R]} properties]{prop}{rsdProp}\label{prop:rsd-properties}
Let $\dbf\in\RSD$. $\dbf$ is element-wise non-negative and $\lone{\dbf}=1$.
\end{restatable}
\begin{proof}
$\dbf$ has non-negative elements because it equals the limit of $\lim_{\gamma\to1}(1-\gamma)\f(\gamma)$, whose elements are non-negative by \cref{prop:visit-dist-prop} \cref{item:mono-increase}.
\begin{align}
    \lone{\dbf}&=\lone{\lim_{\gamma\to1}(1-\gamma)\f(\gamma)}\label{eq:lone-rsd}\\
    &=\lim_{\gamma\to 1} (1-\gamma)\lone{\f(\gamma)}\label{eq:lone-norm-rsd}\\
    &= 1.\label{eq:lone-norm-rsd-1}
\end{align}
\Cref{eq:lone-rsd} follows because the definition of {\rsd}s (\cref{def:rsd}) ensures that $\exists \f\in\F(s):\lim_{\gamma\to1}(1-\gamma)\f(\gamma)=\dbf$. \Cref{eq:lone-norm-rsd} follows because $\lone{\cdot}$ is a continuous function. \Cref{eq:lone-norm-rsd-1} follows because $\lone{\f(\gamma)}=\geom$ by \cref{prop:visit-dist-prop} \cref{item:lone-visit}.
\end{proof}

\begin{restatable}[When reachable with probability 1, {\stateEnd}s induce non-dominated {\rsd}s]{lem}{termNDrsd}\label{lem:term-nd-rsd}
If $\unitvec[s']\in\RSD$, then $\unitvec[s']\in\RSDnd$.
\end{restatable}
\begin{proof}
If $\dbf\in\RSD$ is distinct from $\unitvec[s']$, then $\lone{\dbf}=1$ and $\dbf$ has non-negative entries by \cref{prop:rsd-properties}. Since $\dbf$ is distinct from $\unitvec[s']$, then its entry for index $s'$ must be strictly less than 1: $\dbf^\top\unitvec[s']<1=\unitvec[s']^\top\unitvec[s']$. Therefore, $\unitvec[s']\in\RSD$ is strictly optimal for the \emph{reward function} $\rf\defeq \unitvec[s']$, and so $\unitvec[s']\in\RSDnd$.
\end{proof}

\avgAvoidTerminal*
\begin{proof}
Suppose $\unitvec[s_x],\unitvec[s']\in\RSD$ are distinct. Let $\phi\defeq (s_x \,\,\, s'),D'\defeq \set{\unitvec[s_x]},D\defeq \RSD\setminus\set{\unitvec[s_x]}$. $\phi\cdot D'=\set{\unitvec[s']}\subseteq \RSD\setminus\set{\unitvec[s_x]}\eqdef D$ since $s_x\neq s'$. $D'\cup D=\RSD$ and $\RSDnd\setminus(D'\cup D)=\RSDnd\setminus\RSD=\emptyset$ trivially have pairwise orthogonal vector elements. Then apply \cref{rsdIC} to conclude that $\avgprob[\Dany]{\{\unitvec[s_x]\}}\leqMost[][\DSetAny] \avgprob[\Dany]{\RSD\setminus\{\unitvec[s_x]\}}$.

Suppose there exists another $\unitvec[s'']\in\RSD$. By \cref{lem:term-nd-rsd}, $\unitvec[s'']\in\RSDnd$. Furthermore, since $s''\not \in \set{s',s_x}$ , $\unitvec[s'']\in \prn{\RSD\setminus\set{\unitvec[s_x]}}\setminus \set{\unitvec[s']}=D\setminus \phi\cdot D'$. Therefore, $\unitvec[s'']\in\RSDnd\cap\prn{D\setminus \phi \cdot D'}$. Then apply the second condition of \cref{rsdIC} to conclude that $\avgprob[\Dany]{\{\unitvec[s_x]\}}\not\geqMost[][\DSetBd] \avgprob[\Dany]{\RSD\setminus\{\unitvec[s_x]\}}$. 
\end{proof}